\documentclass[11pt,reqno]{amsart}

\usepackage[T1]{fontenc}
\usepackage{lmodern}
\usepackage[margin=1in]{geometry}
\usepackage{microtype}
\usepackage{amsmath,amssymb,mathtools,mathrsfs}
\usepackage{aliascnt}
\usepackage{booktabs,longtable,array}
\usepackage{tikz-cd}
\usepackage{xcolor}
\usepackage{hyperref}
\usepackage[nameinlink,capitalize,noabbrev]{cleveref}

\definecolor{linknavy}{HTML}{1F3A5F}
\hypersetup{
  colorlinks=true,
  linkcolor=linknavy,
  citecolor=linknavy,
  urlcolor=linknavy,
  pdftitle={Eidetic Power-Map Algebras: Balanced Kernels,
    Rational Classes, and Ramified Completion},
  pdfauthor={Simon Velez}
}

\numberwithin{equation}{section}

\newtheorem{theorem}{Theorem}[section]
\newaliascnt{proposition}{theorem}
\newtheorem{proposition}[proposition]{Proposition}
\aliascntresetthe{proposition}
\newaliascnt{lemma}{theorem}
\newtheorem{lemma}[lemma]{Lemma}
\aliascntresetthe{lemma}
\newaliascnt{corollary}{theorem}
\newtheorem{corollary}[corollary]{Corollary}
\aliascntresetthe{corollary}
\newaliascnt{definition}{theorem}
\newtheorem{definition}[definition]{Definition}
\aliascntresetthe{definition}
\newaliascnt{example}{theorem}
\newtheorem{example}[example]{Example}
\aliascntresetthe{example}
\newaliascnt{problem}{theorem}
\newtheorem{problem}[problem]{Problem}
\aliascntresetthe{problem}
\newaliascnt{conjecture}{theorem}
\newtheorem{conjecture}[conjecture]{Conjecture}
\aliascntresetthe{conjecture}
\crefname{example}{Example}{Examples}
\crefname{problem}{Problem}{Problems}
\crefname{conjecture}{Conjecture}{Conjectures}
\theoremstyle{remark}
\newaliascnt{remark}{theorem}
\newtheorem{remark}[remark]{Remark}
\aliascntresetthe{remark}
\theoremstyle{plain}
\newtheorem{thmintro}{Theorem}
\renewcommand{\thethmintro}{\Alph{thmintro}}
\crefname{thmintro}{Theorem}{Theorems}
\crefname{theorem}{Theorem}{Theorems}
\crefname{proposition}{Proposition}{Propositions}
\crefname{lemma}{Lemma}{Lemmas}
\crefname{corollary}{Corollary}{Corollaries}
\crefname{definition}{Definition}{Definitions}
\crefname{remark}{Remark}{Remarks}

\newcommand{\Q}{\mathbb Q}
\newcommand{\Z}{\mathbb Z}
\newcommand{\C}{\mathbb C}
\newcommand{\M}{\mathbb M}
\newcommand{\B}{\mathbb B}
\newcommand{\N}{\mathbb N}
% Forward semigroup algebra
\newcommand{\Sfwd}{\mathcal S}
\newcommand{\SfwdE}{\mathcal S_E}
% Underlying rational algebra
\newcommand{\PE}{\mathcal P_E}
\newcommand{\PF}{\mathcal P_F}
\newcommand{\Pram}{\mathcal P_{\mathrm{ram}}}
\newcommand{\PAd}{\mathcal P_{\mathrm{Ad}}}
\newcommand{\Pgen}{\mathcal P_{E}}
% Full marked eidetic object
\newcommand{\EidObj}[1]{\mathbf E_{#1}}
\newcommand{\EramObj}{\mathbf E_{\mathrm{ram}}}
\newcommand{\PAdObj}{\mathbf E_{\mathrm{Ad}}}
\newcommand{\End}{\operatorname{End}}
\newcommand{\Alg}{\operatorname{Alg}}
\newcommand{\rank}{\operatorname{rank}}
\newcommand{\Span}{\operatorname{span}}
\newcommand{\Bal}{\operatorname{Bal}}
\newcommand{\supp}{\operatorname{supp}}
\newcommand{\Conj}{\operatorname{Conj}}
\newcommand{\Irr}{\operatorname{Irr}}
\newcommand{\Tr}{\operatorname{Tr}}
\newcommand{\1}{\mathbf 1}

\title[Eidetic power-map algebras]{Eidetic power-map algebras:
balanced kernels, rational classes, and ramified completion}
\author{Simon Velez}
\date{16 September 2026}

\subjclass[2020]{Primary 20C15, 16S50; Secondary 11R11, 11R18, 11R42, 16K20}
\keywords{Adams operations, power maps, conjugacy classes, class functions,
rational semisimple algebras, commutant, balanced kernels, Wedderburn
decomposition, Schur index, \'etale algebras, Dedekind zeta functions,
sporadic simple groups, quadratic Dirichlet characters, cyclotomic fields}

\begin{document}

\begin{abstract}
For a finite group \(G\), let \(V_G\) be the rational class functions with
the inner product weighted by class sizes, and let \(\Psi_n\) be the power
map \(f(C)\mapsto f(C^n)\).  The ramified power-map algebra \(\Pram(G)\)
is the rational algebra generated by \(\Psi_p\) and its adjoint
\(\Psi_p^*\) for the primes \(p\) dividing \(\exp(G)\).  It is an
invariant of the class table enriched by power maps, and it detects
information the ordinary character table does not: \(D_8\) and \(Q_8\)
have dimensions \(17\) and \(10\).

We prove that \(\Pram(G)\) is semisimple and equal to its own bicommutant,
identify the commutant with the balanced kernels on the weighted power
graph, and show that the transfer identity \(\Psi_p^*\Psi_p=M_{r_p}\)
makes the algebra full exactly when the only balanced kernels are scalar.
The power maps prime to the exponent act as Frobenius elements: they
commute with \(\Pram(G)\), they identify \(V_G\) with the Galois
representation of the \'etale algebra of character-value fields, and their
determinants give the complete Euler product of its Dedekind zeta function,
with the ramified factors recovered from inertia invariants.  When the
commutant of \(\Pram(G)\) is commutative, every simple factor carries a
Dirichlet character of the unit group, trivial or quadratic when the center
is \(\Q\), and \(\Pram(G)\) coincides with the algebra generated by all
power maps.

We compute the algebra for groups of prime exponent, cyclic prime-power
groups, dihedral groups of order \(2p\), and all twenty-six sporadic simple
groups.  Every sporadic factor has Schur index one; the centers are \(\Q\),
nine copies of \(\Q(\zeta_3)\), and one \(\Q(\zeta_5)\), and the
nonrational centers occur exactly for the six pariah groups.  For the
Monster, a \(186\)-vector certificate proves
\[
 \Pram(\M)\cong M_{170}(\Q)\oplus M_5(\Q)\oplus M_3(\Q)
 \oplus M_2(\Q)^{\oplus4}\oplus\Q^{\oplus8},
\]
the centralizer of a symmetry group \((C_2)^{14}\) of the weighted power
system.  One generator of that group swaps \(16B\) with \(16C\) and
\(32A\) with \(32B\); it preserves all class sizes and power maps but is
not an automorphism of the character table.
\end{abstract}

\maketitle

\section{Introduction}
\label{sec:introduction}

Let \(G\) be a finite group, \(\Conj(G)\) its set of conjugacy classes
ordered so that the identity class is first, and
\(V_G=\operatorname{Map}(\Conj(G),\Q)\) the space of \(\Q\)-valued class
functions.
For \(n\geq1\) the power map \([x]\mapsto[x^n]\) induces an operator
\[
 (\Psi_n f)(C)=f(C^n),
\]
where \(C^n\) is the class of \(n\)th powers of the elements of \(C\).
The normalized class form
\[
 \langle f,h\rangle_G=\frac1{|G|}\sum_C|C|f(C)h(C)
\]
is positive definite, and its adjoint
\(\Psi_n^*\) records the weighted reverse fibres of the power map.

The construction depends only on the conjugacy classes, their sizes, and
the selected power maps.  It is therefore an invariant of the
power-map-enriched class table.

For a label set \(E\subseteq\N_{>0}\) containing \(1\), write
\begin{equation}
 \PE(G)=\Alg_\Q\{I,\Psi_n,\Psi_n^*:n\in E\}\subseteq\End_\Q(V_G)
 \label{eq:pe-definition}
\end{equation}
for the rational adjoint-closed algebra generated by these operators.
The \(E\)-marked eidetic object is
\begin{equation}
 \EidObj{E}(G)=
 \bigl(V_G,\mathcal B_G,\langle\ ,\ \rangle_G,
       \{\Psi_n\}_{n\in E},\PE(G)\bigr),
 \label{eq:eidetic-object}
\end{equation}
where \(\mathcal B_G\) is the class-coordinate basis.
The canonical label set is the set of primes dividing the exponent,
\(E=\pi(\exp(G))\); then
\(\Pram(G)=\PE(G)\) is the \emph{ramified power-map algebra}.  The
\emph{ramified eidetic object} is
\begin{equation}
 \EramObj(G)=
 \bigl(V_G,\mathcal B_G,\langle\ ,\ \rangle_G,
       \{\Psi_n\}_{n\geq1},\Pram(G)\bigr).
 \label{eq:ramified-eidetic-object}
\end{equation}
The specialization \(E=\N_{>0}\) gives the \emph{full Adams algebra}
\(\PAd(G)\) and its associated object \(\PAdObj(G)\).

\subsection{What the algebra sees}

The algebra uses class sizes and power maps, so it need not factor
through the ordinary complex character table.  It does not.  The smallest
classical example already gives a sharp separation.

\begin{theorem}
\label{thm:d8-q8}
The dihedral group \(D_8\) and quaternion group \(Q_8\) are isoclinic and
have equivalent ordinary complex character tables.  Nevertheless
\[
 \dim_{\Q}\Pram(D_8)=17,\qquad
 \dim_{\Q}\Pram(Q_8)=10,
\]
so their represented rational power-map algebras are nonisomorphic.  Over
\(\C\),
\[
 \Pram(D_8)\otimes_{\Q}\C
 \cong M_4(\C)\oplus\C,
\qquad
 \Pram(Q_8)\otimes_{\Q}\C
 \cong M_3(\C)\oplus\C,
\]
and the scalar \(Q_8\)-constituent has multiplicity two in class-coordinate
space.
\end{theorem}

\begin{proof}
Both groups have five ordinary conjugacy classes of sizes
\(1,1,2,2,2\).  After a permutation of the final three columns, their
common ordinary character table is
\[
\begin{pmatrix}
1& 1& 1& 1& 1\\
1& 1& 1&-1&-1\\
1& 1&-1& 1&-1\\
1& 1&-1&-1& 1\\
2&-2& 0& 0& 0
\end{pmatrix}.
\]
Let \(\rho\) be the two-dimensional irreducible representation.  Then
\[
 \lambda^2(\rho)=\det(\rho)\neq1\quad\text{for }D_8,
 \qquad
 \lambda^2(\rho)=1\quad\text{for }Q_8.
\]
The standard representation of \(D_8\) contains reflections and has
nontrivial determinant, while \(Q_8\subseteq\mathrm{SU}(2)\).  Hence the
ordinary-table identification does not preserve the based augmented
\(\lambda\)-ring.

Their power systems differ.  In \(D_8\), every reflection squares to the
identity.  In \(Q_8\), every noncentral class squares to the nontrivial
central element.  Hence the weighted reverse fibres of the squaring map
differ.

Rational word closure of the squaring and fourth-power pullbacks
together with their class-form adjoints gives dimensions \(17\) and \(10\).
The center and represented-commutant dimensions are
\[
\begin{array}{c|ccc}
 &\dim A&\dim Z(A)&\dim A'\\
\hline
D_8&17&2&2\\
Q_8&10&2&5.
\end{array}
\]
The dimension identities for a semisimple complex algebra acting on a
five-dimensional module have unique solutions
\[
 (d,m)=(4,1),(1,1)
 \quad\text{and}\quad
 (d,m)=(3,1),(1,2),
\]
respectively.  This proves the decompositions and the separation.
\end{proof}

\begin{corollary}
\label{cor:strict-refinement}
The represented power-map algebra is not a function of the ordinary
complex character table, and it is not an isoclinism invariant.
\end{corollary}

The same separation persists after adjoining class sizes and element
orders.  At order \(64\) the first six pairs of groups with matching
character tables, class sizes, and element orders appear; the signature
\((\dim A,\dim Z(A),\dim A')\) separates two of them, so the
power-map algebra does not factor through these data.

\subsection{The Monster}

The largest case we compute is the Monster \(\M\), with \(194\) classes
and exponent divisible by fifteen primes.  Let \(\Gamma_{\M}\) be the group
of permutations of \(\Conj(\M)\) that preserve class sizes and commute
with every prime-power map; it is the automorphism group of the weighted
power system.

\begin{thmintro}[\cref{thm:monster-symmetry-commutant}]
\label{thm:intro-monster}
\(\Gamma_{\M}\) is elementary abelian of rank \(14\), and
\[
 \Pram(\M)
 =\End_{\Q\Gamma_{\M}}(V_{\M})
 \cong
 M_{170}(\Q)\oplus M_5(\Q)\oplus M_3(\Q)
 \oplus M_2(\Q)^{\oplus4}\oplus\Q^{\oplus8},
\]
of dimension \(28958\), with center and commutant of dimension \(15\).
The subgroup of \(\Gamma_{\M}\) induced by Galois conjugation of classes
has rank \(13\).  A complement is generated by the involution
\[
 (16B\ 16C)(32A\ 32B),
\]
which preserves every class size and every power map but is not an
automorphism of the character table.
\end{thmintro}

The equality \(\Pram(\M)=\End_{\Q\Gamma_{\M}}(V_{\M})\) says that the
algebra is as large as the symmetry group allows: the weighted power graph
of the Monster has no fractional symmetries beyond its automorphisms.  The
proof is a saturation certificate of \(186\) orbit vectors
(\cref{sec:rank-one-saturation}).  The complementary involution is a
symmetry of the power-map-enriched class table invisible to the character
table.  It accounts for the one nonprincipal block of \(\Pram(\M)\) on
which Galois conjugation acts trivially.

\subsection{Main results}

The general theory behind \cref{thm:intro-monster} is the following.  Write
\(N=\exp(G)\) and \(U_N=(\Z/N)^\times\).  The results are stated for the
group case; \cref{thm:intro-structure} holds for any finite weighted power
system (\cref{def:weighted-power-system}).

\begin{thmintro}[Structure; \cref{thm:rational-bicommutant-structural,thm:balanced-commutant,cor:full-matrix}]
\label{thm:intro-structure}
\(\PE(G)\) is a semisimple \(\Q\)-algebra equal to its own bicommutant in
\(\End_\Q(V_G)\).  Its commutant is the space \(\Bal_E(G)\) of balanced
kernels on the weighted power graph, so \(\PE(G)=\Bal_E(G)'\), and
\(\PE(G)=\End_\Q(V_G)\) if and only if the only balanced kernels are the
scalar multiples of the identity.  For \(G\neq1\) the algebra
\(\Pram(G)\) is noncommutative.
\end{thmintro}

\begin{thmintro}[Frobenius; \cref{prop:good-in-commutant,thm:rational-class-decomposition,cor:good-euler-product,thm:ramified-completion,thm:good-label-characters,thm:good-label-saturation}]
\label{thm:intro-frobenius}
Let \(u\in U_N\).  Then \(\Psi_u\) is an isometry of \(V_G\) commuting
with \(\Pram(G)\).  As a \(\Q[U_N]\)-module,
\[
 V_G\cong\mathcal E(G)=\prod_{R\in\Conj(G)/U_N}F_R,
\]
where \(F_R\) is the field generated by the values of the irreducible
characters on the classes in the rational class \(R\).  For \(p\nmid N\)
the operator \(\Psi_p\) is the Frobenius at \(p\), and
\[
 \prod_{p\nmid N}\det(I-\Psi_pp^{-s})^{-1}
 =\zeta^{(N)}_{\mathcal E(G)}(s).
\]
The factors at \(p\mid N\) are recovered from inertia invariants of the
good-label operators on each element-order layer, giving the complete
Euler product of \(\zeta_{\mathcal E(G)}\).  If the commutant of
\(\Pram(G)\) is commutative, every simple factor \(A_k\) with center
\(K_k\) carries a character \(\lambda_k:U_N\to\mu(K_k)\), trivial or
quadratic when \(K_k=\Q\), and \(\Pram(G)=\PAd(G)\).
\end{thmintro}

\begin{thmintro}[Census; \cref{thm:sporadic-census,thm:good-label-census,cor:nonrational-from-orbits}]
\label{thm:intro-census}
For the twenty-six sporadic simple groups, \(\Pram(G)\) is determined
explicitly.  Every simple factor has Schur index one.  The commutant is
commutative for twenty-two groups, the exceptions being \(O'N\),
\(Fi_{24}'\), \(HN\), and \(\B\).  The center fields are \(\Q\), nine
copies of \(\Q(\zeta_3)\), and one copy of \(\Q(\zeta_5)\); the
nonrational centers occur exactly for the six pariah groups
\(J_1,J_3,J_4,Ru,O'N,Ly\), and each comes from a single rational class of
prime size three or five.
\end{thmintro}

\subsection{Related constructions}

Element-level functional graphs of the power map \(x\mapsto x^n\) retain
the cycle and rooted-tree structure on all elements, as studied for finite
abelian and several nonabelian families by Qureshi and
Reis~\cite{QureshiReis}.  The present construction retains one vertex per
conjugacy class, equips the quotient graph with the nonuniform class
measure, and uses several labels simultaneously; the class-form adjoint
records every reverse fibre weighted by class size, so two class-power
graphs with the same unweighted shape can produce different algebras.
The forward operators alone are the linearization of the transformation
monoid generated by the \(\tau_n\)~\cite{SteinbergTransformation,
SteinbergMonoids}; the algebra studied here is the star closure of that
image, and it is contained in the coherent closure of the colored
digraph~\cite{HigmanCoherent,HigmanCoherentAlgebras,Weisfeiler,BannaiIto,
ChenPonomarenko} whenever that closure contains the class-measure diagonal
and its inverse.  GAP's documentation notes that power maps are not in
general determined by the matrix of irreducible characters
alone~\cite{GAPPowerMaps}; \cref{thm:d8-q8} is the algebraic form of that
remark.

\subsection{Outline}

\Cref{sec:algebra} defines finite weighted power systems, proves the
transfer calculus and \cref{thm:intro-structure}, and derives the
saturation and degeneracy criteria.  \Cref{sec:goodlabels} treats the
Galois action and proves \cref{thm:intro-frobenius}.  \Cref{sec:sporadic}
contains the sporadic census, the good-label characters, the rank-one
saturation method, and the proof of \cref{thm:intro-monster}, together
with the symmetries invisible to the character table.
\Cref{sec:families} computes the algebra for several infinite families and
records a comparison panel of small groups.  \Cref{sec:certification}
describes the computations, and \cref{sec:open-problems} collects open
questions.  \Cref{sec:power-dynamics} records a dynamical invariant of the
prime-power maps that separates the sporadic groups.

\section{The algebra and its commutant}
\label{sec:algebra}

\subsection{Finite weighted power systems}
\label{sec:systems}

The finite weighted dynamics that the algebra linearizes is the following.

\begin{definition}
\label{def:weighted-power-system}
Fix a finite label set \(E\).  A \emph{finite weighted power system} is a
tuple
\[
 \mathscr X=(X,\mu,x_0,(\tau_e)_{e\in E}),
\]
where \(X\) is a finite set, \(\mu:X\to\Q_{>0}\) has total mass one,
\(x_0\in X\), and the maps \(\tau_e:X\to X\) commute and fix \(x_0\).
An isomorphism \(a:\mathscr X\to\mathscr Y\) is a bijection satisfying
\[
 a(x_0)=y_0,\qquad \nu(a(x))=\mu(x),\qquad
 a\tau_e=\upsilon_ea\quad(e\in E).
\]
\end{definition}

\begin{definition}
\label{def:multiplicative-labels}
Call the system \emph{multiplicatively labeled} if \(E\) is a set of positive
integers, \(1\in E\), and
\begin{equation}
 \tau_m\tau_n=\tau_{mn}\ \text{whenever }m,n,mn\in E,
 \qquad
 \tau_1=\mathrm{id}.
 \label{eq:label-multiplicativity}
\end{equation}
This hypothesis is \emph{not} part of \cref{def:weighted-power-system}; it is
invoked explicitly where it is used, namely in
\cref{prop:minimal-prime-generators} and throughout
\cref{sec:goodlabels}.  The structural results of the present section do
not require it.
\end{definition}

Write \(V_X=\operatorname{Map}(X,\Q)\), let
\(\mathcal B_X=(\delta_x)_{x\in X}\) be its coordinate basis, and put
\[
 \langle f,h\rangle_\mu=\sum_{x\in X}\mu(x)f(x)h(x).
\]
The pullback \(P_e f=f\circ\tau_e\) has coordinate matrix
\begin{equation}
 (P_e)_{xy}=\mathbf1_{\{y=\tau_e(x)\}},
 \qquad
 P_e^*=D_\mu^{-1}P_e^{\mathsf T}D_\mu,
 \quad D_\mu=\operatorname{diag}(\mu(x)).
 \label{eq:system-adjoint}
\end{equation}
Let \(\mathbf c_\mu=(c_x)_{x\in X}\) be the primitive positive integral
vector proportional to \(\mu\), and regard both
\(\mathbf c_\mu=\sum_xc_x\delta_x\) and
\(\eta_0=\delta_{x_0}\) as marked vectors.

\begin{definition}
\label{def:fully-marked-object}
The fully marked linear object attached to \(\mathscr X\) is
\[
 \mathbf P(\mathscr X)=
 \bigl(
 V_X,\mathcal B_X,\langle\ ,\ \rangle_\mu,\mathbf1,\eta_0,
 \mathbf c_\mu,(P_e)_{e\in E},
 \PE(\mathscr X)
 \bigr),
\]
where
\[
 \PE(\mathscr X)
 =\Alg_\Q\{I,P_e,P_e^*:e\in E\}\subseteq\End_\Q(V_X).
\]
\end{definition}

The forward subalgebra
\(\SfwdE(\mathscr X)=\Alg_\Q\{I,P_e:e\in E\}\) is commutative, since the
\(\tau_e\) commute; \(\PE(\mathscr X)\) is its star closure.
If \(a:\mathscr X\to\mathscr Y\) is an isomorphism, then
\((U_af)(y)=f(a^{-1}y)\) carries every part of the marking to the
corresponding part, is an isometry, and intertwines \(P_e^X\) with
\(P_e^Y\) and their adjoints.  Conversely an isomorphism of fully marked
objects permutes their coordinate bases and therefore comes from a unique
system isomorphism.  The fully marked object thus reconstructs
\((X,\mu,x_0,(\tau_e))\) up to unique simultaneous relabeling.

More general factor maps do not act on the star closure.  Suppose
\(q:(X,\mu,\tau_e)\to(Y,\nu,\upsilon_e)\) is a surjection with
\(q_*\mu=\nu\) and \(q\tau_e=\upsilon_eq\).  Pullback
\(U_q:V_Y\to V_X\), \(U_qf=f\circ q\), is an isometry and
\begin{equation}
 P_e^XU_q=U_qP_e^Y.
 \label{eq:factor-forward-intertwining}
\end{equation}
It follows only that
\(U_q^*(P_e^X)^*=(P_e^Y)^*U_q^*\).  The stronger relation
\((P_e^X)^*U_q=U_q(P_e^Y)^*\), needed to intertwine the generated
adjoint-closed algebras, holds if and only if
\(\operatorname{im}U_q\) is reducing for \(P_e^X\).  Equivalently, for
every \(x\in X\) and \(z\in Y\),
\begin{equation}
 \frac1{\mu(x)}
 \sum_{\substack{\tau_e y=x\\q(y)=z}}\mu(y)
 =
 \begin{cases}
 \displaystyle\frac{\nu(z)}{\nu(q(x))},
   &\upsilon_ez=q(x),\\[6pt]
 0,&\upsilon_ez\ne q(x).
 \end{cases}
 \label{eq:reverse-fibre-condition}
\end{equation}
General measure-preserving factor maps need not satisfy this
reverse-fibre condition.

\subsubsection{The system supplied by a finite group}

Let
\(\operatorname{Conj}(G)=\{C_1,\ldots,C_k\}\), with
\(C_1=\{1\}\), and set
\begin{equation}
 V_G=\operatorname{Map}(\operatorname{Conj}(G),\Q),\qquad
 \langle f,h\rangle_G
 =\frac1{|G|}\sum_{i=1}^k |C_i|f(C_i)h(C_i).
 \label{eq:coordinate-class-form}
\end{equation}
For \(n\geq1\), write \(\sigma_n(i)=j\) when \(g_i^n\in C_j\).
Then
\begin{equation}
 (P_n)_{ij}=\mathbf1_{\{j=\sigma_n(i)\}},
 \qquad
 P_n^*=D_G^{-1}P_n^{\mathsf T}D_G,\quad
 D_G=\operatorname{diag}\left(\frac{|C_i|}{|G|}\right).
 \label{eq:power-matrix}
\end{equation}
Thus \(G\) supplies the finite weighted power system
\[
 \mathscr C_E(G)=
 \left(\operatorname{Conj}(G),\frac{|C|}{|G|},\{1\},
 ([x]\mapsto[x^n])_{n\in E}\right).
\]
Its primitive weight vector is the positive class-size vector
\((|C_i|)_i\).  We define
\begin{equation}
 \mathbf P(G,E)
 =\mathbf P(\mathscr C_E(G)),\qquad
 \PE(G)
 =\PE(\mathscr C_E(G)).
 \label{eq:ap-definition-structural}
\end{equation}
For \(N=\exp(G)\), the canonical label set is
\(E=\pi(N)\), and we omit \(E\) from the notation.

A group isomorphism induces an isomorphism of these fully marked objects.
More generally, simultaneous relabeling of classes preserving class sizes,
the identity class, and all selected power maps is isomorphism at
the source-system level.

\begin{proposition}
\label{prop:quotient-factor}
If \(q:G\twoheadrightarrow Q\) is a quotient, then
\[
 \bar q:\operatorname{Conj}(G)\longrightarrow\operatorname{Conj}(Q),
 \qquad [g]\longmapsto[q(g)],
\]
is a surjective measure-preserving factor map for every power map.
Consequently pullback \(U_{\bar q}:V_Q\to V_G\) is an isometry and
\(\Psi_n^GU_{\bar q}=U_{\bar q}\Psi_n^Q\).
It intertwines the full algebras if and only if its image
is reducing for every \(\Psi_n^G\).
\end{proposition}

\begin{proof}
The map on classes is well defined, surjective, and commutes with powers.
For a class \(D\subseteq Q\), its full inverse image has
\(|\ker q||D|\) elements.  Hence the sum of the normalized \(G\)-class
weights above \(D\) is
\[
 \frac{|\ker q||D|}{|G|}=\frac{|D|}{|Q|}.
\]
The remaining assertions are \eqref{eq:factor-forward-intertwining} and its
reverse-fibre criterion.
\end{proof}

A proper subgroup inclusion does not preserve normalized class measure
(the identity class has mass \(1/|H|\) in \(H\) and \(1/|G|\) in \(G\)),
so no proper inclusion is a morphism of weighted power systems.

\subsection{Transfer operators}
\label{sec:transfer}

Let \(\mathscr X=(X,\mu,x_0,(\tau_e)_{e\in E})\) be a finite weighted
power system.  Write
\[
 P_ef=f\circ\tau_e,\qquad L_e=P_e^*,
\]
and let \(M_a\) denote multiplication by \(a\in V_X\).  The reverse-fibre
density is
\begin{equation}
 r_e(y)=\frac{1}{\mu(y)}
 \sum_{\tau_ex=y}\mu(x).
 \label{eq:reverse-fibre-density}
\end{equation}

\begin{theorem}
\label{thm:transfer-calculus}
For \(a,b\in V_X\) and \(e,f\in E\),
\begin{align}
 P_eM_a&=M_{P_ea}P_e,                         \label{eq:covariance}\\
 L_e\bigl((P_ea)b\bigr)&=aL_eb,               \label{eq:frobenius-transfer}\\
 L_eM_bP_e&=M_{L_eb},                         \label{eq:transfer-conjugation}\\
 L_eP_e&=M_{r_e}.                             \label{eq:gram-density}
\end{align}
If \(\tau_{ef}=\tau_e\tau_f=\tau_f\tau_e\), then
\begin{equation}
 L_eL_f=L_fL_e=L_{ef},\qquad
 r_{ef}=L_er_f=L_fr_e.
 \label{eq:density-cocycle}
\end{equation}
\end{theorem}

\begin{proof}
The first identity follows from
\[
 (P_e(ah))(x)=a(\tau_ex)h(\tau_ex).
\]
The adjoint is
\begin{equation}
 (L_eb)(y)=\frac{1}{\mu(y)}
 \sum_{\tau_ex=y}\mu(x)b(x).
 \label{eq:explicit-transfer}
\end{equation}
On the fibre \(\tau_ex=y\), the function \(P_ea\) has value \(a(y)\).
This gives \eqref{eq:frobenius-transfer} and
\eqref{eq:transfer-conjugation}.  Taking \(b=1\) gives
\eqref{eq:gram-density}.  The last two identities follow by taking
adjoints in \(P_eP_f=P_fP_e=P_{ef}\) and applying the result to \(1\).
\end{proof}

\begin{theorem}
\label{thm:normality-commutativity}
For a label \(e\), the following conditions are equivalent:
\begin{enumerate}
\item \(P_e\) is normal;
\item \(P_e\) is unitary;
\item \(\tau_e\) is a measure-preserving permutation;
\item \(r_e=1\) and \(\tau_e\) is injective.
\end{enumerate}
Consequently
\begin{equation}
 \PE(\mathscr X)\text{ is commutative}
 \quad\Longleftrightarrow\quad
 \text{every \(\tau_e\) is a measure-preserving permutation.}
 \label{eq:commutative-dichotomy}
\end{equation}
\end{theorem}

\begin{proof}
The matrix \(L_eP_e=M_{r_e}\) is diagonal, while
\begin{equation}
 (P_eL_e)_{xz}=
 \begin{cases}
 \mu(z)/\mu(\tau_ex),&\tau_ez=\tau_ex,\\
 0,&\tau_ez\neq\tau_ex.
 \end{cases}
 \label{eq:fibre-gram}
\end{equation}
If \(P_e\) is normal, \eqref{eq:fibre-gram} is diagonal.  Every fibre of
\(\tau_e\) then has at most one point, so \(\tau_e\) is a permutation.
For a permutation,
\[
 r_e(y)=\frac{\mu(\tau_e^{-1}y)}{\mu(y)}.
\]
Normality gives \(r_e(x)=r_e(\tau_ex)\).  This value is constant on each
cycle, and its product around the cycle is \(1\); hence \(r_e=1\).
Thus \(\tau_e\) preserves \(\mu\) and \(P_e\) is unitary.  The remaining
implications are immediate.  Commuting measure-preserving permutations
have commuting inverse pullbacks, which proves
\eqref{eq:commutative-dichotomy}.
\end{proof}

\begin{corollary}
\label{cor:group-normal-labels}
Let \(G\) be finite and \(N=\exp(G)\).  Then
\begin{equation}
 \Psi_n\text{ is normal}
 \quad\Longleftrightarrow\quad
 \Psi_n\text{ is unitary}
 \quad\Longleftrightarrow\quad
 (n,N)=1.
 \label{eq:good-normal-unitary}
\end{equation}
The good labels generate a commutative semisimple algebra.  Every prime
\(p\mid N\) gives a nonnormal generator; hence
\(\Pram(G)\) is noncommutative for \(G\neq1\).
\end{corollary}

\begin{proof}
If \((n,N)=1\), choose \(m\) with \(mn\equiv1\pmod N\).  The \(n\)-th and
\(m\)-th power maps are inverse permutations of the conjugacy classes and
preserve their sizes.  If \((n,N)>1\), choose
\(p\mid(n,N)\) and an element \(x\) of order \(p\).  Then \(x^n=1\), so
the power map on conjugacy classes is not injective.  Apply
\cref{thm:normality-commutativity}.
\end{proof}

For a finite group, \eqref{eq:reverse-fibre-density} becomes
\begin{equation}
 r_n(C)=
 \frac{\#\{x\in G:x^n\in C\}}{|C|},
 \qquad
 r_n(\{1\})=\#\{x\in G:x^n=1\}.
 \label{eq:group-ramification-density}
\end{equation}
Thus \(\Psi_n^*\Psi_n\) is the diagonal operator of class-resolved
\(n\)-th root counts.

\begin{definition}
\label{def:ramification-algebra}
Let \(\mathcal R_E(\mathscr X)\) be the smallest unital
\(\Q\)-subalgebra of \(V_X\) stable under every \(L_e\), and put
\[
 \mathcal D_E(\mathscr X)
 =\{M_a:a\in\mathcal R_E(\mathscr X)\}.
\]
\end{definition}

\begin{proposition}
\label{prop:ramification-diagonal-inclusion}
For every finite weighted power system,
\[
 \mathcal D_E(\mathscr X)\subseteq\PE(\mathscr X).
\]
\end{proposition}

\begin{proof}
The identity multiplier belongs to \(\PE(\mathscr X)\).  If \(M_a\)
belongs to it, then
\[
 M_{L_ea}=L_eM_aP_e\in\PE(\mathscr X)
\]
by \eqref{eq:transfer-conjugation}.  The multipliers are closed under sums
and products.
\end{proof}

\begin{theorem}
\label{thm:diagonal-separation}
Let \(\Gamma_E\) be the undirected graph on \(X\) with edges
\(\{x,\tau_ex\}\).  If \(\mathcal R_E(\mathscr X)\) separates the points
of \(X\), then
\begin{equation}
 \PE(\mathscr X)
 \cong\bigoplus_{C\in\pi_0(\Gamma_E)}
       \End_\Q(\Q^C).
 \label{eq:component-matrix-decomposition}
\end{equation}
If \(\Gamma_E\) is connected, then
\(\PE(\mathscr X)=\End_\Q(V_X)\).
\end{theorem}

\begin{proof}
For \(x\in X\) and \(y\neq x\), choose \(a_y\in\mathcal R_E(\mathscr X)\)
with \(a_y(x)\neq a_y(y)\).  Then
\[
 q_x=\prod_{y\neq x}
 \frac{a_y-a_y(y)}{a_y(x)-a_y(y)}
 =\mathbf1_{\{x\}}.
\]
Hence all coordinate idempotents \(E_{xx}=M_{\mathbf1_{\{x\}}}\) belong
to \(\PE(\mathscr X)\).  If \(\tau_ex=y\), then
\[
 E_{xx}P_eE_{yy}=E_{xy},\qquad
 E_{yy}L_eE_{xx}=\frac{\mu(x)}{\mu(y)}E_{yx}.
\]
Products along paths in \(\Gamma_E\) give every matrix unit within each
connected component.  Every generator preserves the direct sum over the
components.
\end{proof}

\subsection{Semisimplicity and the rational bicommutant}
\label{sec:structural-theory}

\begin{proposition}
\label{prop:minimal-prime-generators}
Let \(N=\exp(G)\).  Each of the following sets generates the same unital
adjoint-closed rational algebra:
\begin{equation}
 \{\Psi_p:p\mid N,\ p\text{ prime}\},\qquad
 \{\Psi_{p^a}:p^a\mid N\},\qquad
 \{\Psi_d:d\mid N\}.
 \label{eq:equivalent-exponent-sets}
\end{equation}
Here the adjoints of the displayed operators are included in every case.
This common algebra is \(\Pram(G)\).
\end{proposition}

\begin{proof}
The identities
\[
 \Psi_m\Psi_n=\Psi_{mn},\qquad \Psi_1=I
\]
show that \(\Psi_{p^a}=\Psi_p^a\) and, for
\(d=\prod_p p^{a_p}\), that \(\Psi_d=\prod_p\Psi_p^{a_p}\).
Thus all divisor maps lie in the algebra generated by the prime maps.  The
reverse inclusion is immediate.  Since an adjoint-closed algebra contains
the adjoint of every product it contains, adjoining the listed adjoints
does not alter this conclusion.
\end{proof}

\begin{example}
\label{ex:all-positive-exponents}
The algebra generated by \(\Psi_n\) for every \(n\geq1\) need not equal
\(\Pram(G)\).  For example, if \(G=C_5\), the canonical
algebra uses only \(\Psi_5\), whereas \(\Psi_2\) cyclically permutes the
four nonidentity class coordinates.  The prime-exponent calculation in
\cref{thm:prime-exponent-structural} shows that the canonical algebra acts
scalarly on a three-dimensional complementary summand, while \(\Psi_2\)
does not.  Hence the all-exponents algebra is strictly larger in this
example.  Proposition~\ref{prop:minimal-prime-generators} concerns the
power maps indexed by divisors of \(\exp(G)\), not all residue classes
modulo the exponent.
\end{example}

For a subalgebra \(A\subseteq\End_F(V)\), write \(A'\) for its commutant
in \(\End_F(V)\).

\begin{lemma}
\label{lem:commutant-base-change}
Let \(F\subseteq K\) be a field extension, let \(V\) be a
finite-dimensional \(F\)-vector space, and let
\(A\subseteq\End_F(V)\) be a finite-dimensional subalgebra.  Under the
natural identification
\(\End_F(V)\otimes_FK\cong\End_K(V\otimes_FK)\), one has
\begin{equation}
 A'\otimes_FK=(A\otimes_FK)'.
 \label{eq:commutant-base-change}
\end{equation}
\end{lemma}

\begin{proof}
Choose an \(F\)-basis \(a_1,\ldots,a_r\) of \(A\).  Its commutant is the
kernel of the \(F\)-linear map
\[
 T:\End_F(V)\longrightarrow\End_F(V)^{\oplus r},
 \qquad X\longmapsto([X,a_1],\ldots,[X,a_r]).
\]
Tensoring the exact sequence
\(0\to\ker T\to\End_F(V)\xrightarrow{T}\End_F(V)^{\oplus r}\)
with the field \(K\) preserves exactness.  The resulting kernel is
the commutant of \(A\otimes_FK\).
\end{proof}

\begin{theorem}
\label{thm:rational-bicommutant-structural}
Let \(X\) be a finite set, let \(\mu:X\to\Q_{>0}\), and let
\((T_e)_{e\in E}\) be any finite family of \(\Q\)-linear operators on
\(V_X\), with \(T_e^*=D_\mu^{-1}T_e^{\mathsf T}D_\mu\) the adjoint for
\(\langle\ ,\ \rangle_\mu\).  Then
\[
 A=\Alg_\Q\{I,T_e,T_e^*:e\in E\}\subseteq\End_\Q(V_X)
\]
is a finite-dimensional semisimple \(\Q\)-algebra, and in its defining
faithful representation on \(V_X\),
\begin{equation}
 A''=A.
 \label{eq:rational-double-commutant-structural}
\end{equation}
Over \(\C\), there are pairwise inequivalent simple modules
\(W_1,\ldots,W_t\), with \(d_j=\dim_\C W_j\), and multiplicity spaces
\(M_j\), with \(m_j=\dim_\C M_j\), such that
\begin{equation}
 V_X\otimes_\Q\C\cong\bigoplus_{j=1}^t W_j\otimes_\C M_j,
 \qquad
 A\otimes_\Q\C\cong
 \bigoplus_{j=1}^t\End_\C(W_j)\otimes I_{M_j}.
 \label{eq:complex-isotypic-structural}
\end{equation}
In particular,
\begin{equation}
 \dim_\C(A\otimes\C)=\sum_jd_j^2,
 \qquad
 \dim_\C(A\otimes\C)'=\sum_jm_j^2.
 \label{eq:dimension-identities-structural}
\end{equation}
In particular this applies to \(A=\PE(\mathscr X)\) for any finite weighted
power system, and to \(A=\PE(G)\) for any finite group and any exponent set.
\end{theorem}

\Cref{thm:rational-bicommutant-structural} uses none of the structure of a
finite weighted power system beyond positivity of \(\mu\): the \(T_e\) need
not commute, need not be pullbacks along maps, and no distinguished point is
required.  The star closure alone gives semisimplicity and
bicommutancy.  Deterministic commuting maps are what give the balanced-kernel
description of the commutant (\cref{thm:balanced-commutant}), and
multiplicative labels are what give the good-label arithmetic of
\cref{sec:goodlabels}.

\begin{proof}
The algebra is finite-dimensional because it is a subalgebra of
\(\End_\Q(V_X)\).  Extend scalars and the form \(\langle\ ,\ \rangle_\mu\) to
obtain a Hermitian form on \(V_X\otimes_\Q\C\), positive definite because
\(\mu>0\) pointwise.  The complex
algebra \(A_\C=A\otimes_\Q\C\) is closed under its Hilbert-space adjoint.
Its Jacobson radical \(J\) is preserved by the involution because the
involution is an anti-automorphism and the radical is invariant under
anti-automorphisms.  If \(x\in J\), then \(x^*x\in J\).  The ideal \(J\)
is nilpotent, while \(x^*x\) is positive semidefinite, so \(x^*x=0\) and
therefore \(x=0\).  Hence \(A_\C\), and consequently \(A\), is
semisimple.

The finite-dimensional structure theorem gives the
isotypic form \eqref{eq:complex-isotypic-structural}~\cite{Wedderburn}.
On each summand,
its commutant is
\(I_{W_j}\otimes\End_\C(M_j)\), whose commutant is
\(\End_\C(W_j)\otimes I_{M_j}\).  Thus \(A_\C''=A_\C\).  Applying
\cref{lem:commutant-base-change} twice gives
\[
 A''\otimes_\Q\C=A_\C''=A_\C=A\otimes_\Q\C.
\]
Since \(A\subseteq A''\) and scalar extension preserves dimensions, this
forces \(A=A''\).  The dimension identities follow from the two displayed
block actions.  The final sentence is the case \(T_e=P_e\) of
\eqref{eq:system-adjoint}.
\end{proof}

\begin{proposition}
\label{prop:field-extension}
For every field extension \(K/\Q\),
\[
 \PE(G)\otimes_\Q K
 =\Alg_K\{I,\Psi_n,\Psi_n^*:n\in E\}
 \subseteq\End_K(V_G\otimes_\Q K),
\]
and
\[
 \PE(G)'\otimes_\Q K
 =(\PE(G)\otimes_\Q K)'.
\]
Thus extension of scalars preserves the represented algebra and its
commutant.  The center fields, Galois orbits, and Schur indices give the
descent data that recover the rational form from its complex blocks.
\end{proposition}

\begin{proof}
The first identity follows because the same rational generators span both
sides after scalar extension.  The second is
\cref{lem:commutant-base-change}.  The final statement is the usual
distinction between a semisimple rational algebra and its split complex
form.
\end{proof}

Direct products exhibit a natural tensor relation, but the power
maps on the two factors are synchronized by the same exponent.

\begin{proposition}
\label{prop:direct-product-containment}
For finite groups \(G,H\) and any exponent set \(E\), the canonical
identification
\[
 V_{G\times H}\cong V_G\otimes_\Q V_H
\]
satisfies
\begin{equation}
 \Psi_n^{G\times H}=\Psi_n^G\otimes\Psi_n^H,
 \qquad
 (\Psi_n^{G\times H})^*=(\Psi_n^G)^*\otimes(\Psi_n^H)^*.
 \label{eq:direct-product-generators}
\end{equation}
Consequently,
\begin{equation}
 \PE(G\times H)
 \subseteq
 \PE(G)\otimes_\Q
 \PE(H)
 \label{eq:direct-product-containment}
\end{equation}
as represented algebras.  For canonical exponent sets, take
\(E=\pi(\exp(G\times H))\).  If \(G\) and \(H\) have the same prime
divisors in their exponents, the right-hand side is
\(\Pram(G)\otimes
\Pram(H)\).
\end{proposition}

\begin{proof}
The conjugacy classes of \(G\times H\) are the products \(C_i\times D_j\),
their sizes multiply, and
\((g,h)^n=(g^n,h^n)\).  Hence the power matrix and class-measure matrix are
\[
 P_n^{G\times H}=P_n^G\otimes P_n^H,
 \qquad
 D_{G\times H}=D_G\otimes D_H.
\]
Formula \eqref{eq:system-adjoint} gives the second identity in
\eqref{eq:direct-product-generators}.  Every generator on the left of
\eqref{eq:direct-product-containment} therefore belongs to the tensor
product algebra on the right.
\end{proof}

Containment in \eqref{eq:direct-product-containment} can be strict.  The
left-hand algebra is generated by the synchronized tensors
\(\Psi_n^G\otimes\Psi_n^H\) and their adjoints.  Tensor-product equality
holds if and only if these synchronized generators recover the separate
operators \(\Psi_n^G\otimes I\) and \(I\otimes\Psi_n^H\); strict examples
occur when they do not.

\subsection{Balanced kernels}
\label{sec:reconstruction}

The commutant has an intrinsic description directly on the weighted colored
digraph.  Write an operator \(S\in\End_\Q(V_X)\) in coordinate form
\[
 (Sf)(x)=\sum_{y\in X}s(x,y)f(y).
\]

\begin{definition}
\label{def:balanced-kernels}
Let \(\operatorname{Bal}(\mathscr X)\) be the rational vector space of
kernels \(s:X\times X\to\Q\) satisfying, for every \(e\in E\) and
\(x,z\in X\),
\begin{align}
 s(\tau_e x,z)
 &=\sum_{\tau_e y=z}s(x,y),
 \label{eq:forward-balance}\\
 \frac{\mu(z)}{\mu(\tau_e z)}s(x,\tau_e z)
 &=\frac1{\mu(x)}
   \sum_{\tau_e y=x}\mu(y)s(y,z).
 \label{eq:reverse-balance}
\end{align}
Composition is ordinary matrix multiplication.
\end{definition}

These kernels balance the colored forward fibres and their weighted reverse
fibres.

\begin{theorem}
\label{thm:balanced-commutant}
For every finite weighted power system,
\begin{equation}
 \PE(\mathscr X)'=\operatorname{Bal}(\mathscr X).
 \label{eq:balanced-commutant}
\end{equation}
Consequently
\begin{equation}
 \PE(\mathscr X)
 =\operatorname{Bal}(\mathscr X)'
 \label{eq:balanced-bicommutant}
\end{equation}
in its defining representation.  In particular the algebra is determined
by the fractional symmetries of the weighted colored power graph.
\end{theorem}

\begin{proof}
Matrix multiplication gives
\[
 (SP_e)_{xz}=\sum_{\tau_e y=z}s(x,y),\qquad
 (P_eS)_{xz}=s(\tau_ex,z),
\]
so \(SP_e=P_eS\) is \eqref{eq:forward-balance}.  From
\((P_e^*)_{yz}=\mu(z)\mu(y)^{-1}\mathbf1_{\{y=\tau_ez\}}\), one obtains
\[
 (SP_e^*)_{xz}=\frac{\mu(z)}{\mu(\tau_ez)}s(x,\tau_ez),\qquad
 (P_e^*S)_{xz}=\frac1{\mu(x)}
 \sum_{\tau_ey=x}\mu(y)s(y,z).
\]
Thus \(SP_e^*=P_e^*S\) is \eqref{eq:reverse-balance}.  Commuting with all
generators is equivalent to commuting with their generated algebra, proving
\eqref{eq:balanced-commutant}.  Equation
\eqref{eq:balanced-bicommutant} is the rational double-commutant theorem.
\end{proof}

Every source automorphism gives a permutation kernel in
\(\operatorname{Bal}(\mathscr X)\).  Automorphisms alone usually do not
describe the commutant: arbitrary rational balanced
kernels are allowed.  For a group of prime exponent with \(k\) classes, the
decomposition in \cref{thm:prime-exponent-structural} gives the explicit
formula
\[
 \Pram(G)'
 \cong \Q\oplus M_{k-2}(\Q),\qquad
 \dim_\Q\Pram(G)'=1+(k-2)^2.
\]
This is a family-level commutant calculation rather than a matrix nullspace
computation.

Between the fully marked object and the abstract algebra there are two
further levels of forgetting: the operator-marked object
\((V,\langle\ ,\ \rangle,\mathbf1,\eta_0,(P_e))\), which retains the
labeled generators but not the coordinate basis, and the unmarked
represented algebra \((A\subseteq\End_\Q(V),V,\langle\ ,\ \rangle)\).
Each arrow in this tower is strict over \(\Q\): the last by the
prime-exponent family, where \(C_3\) and \(C_5\) share the abstract algebra
\(M_2(\Q)\oplus\Q\) but have scalar factors of different multiplicity; the
others by explicit small systems, which we omit.

\subsection{Saturation and two degeneracy criteria}
\label{sec:saturation}

The balanced-kernel theorem converts three questions about the algebra
into questions about balanced kernels on the power graph: when enlarging the
label set changes nothing, when the algebra degenerates to everything, and
when the commutant is commutative.  All three are consequences of
\cref{thm:balanced-commutant} together with the rational double commutant.

\begin{theorem}
\label{thm:saturation}
Let \(E\subseteq F\) be finite label sets for \(\mathscr X\).  Then
\(\PE(\mathscr X)\subseteq\PF(\mathscr X)\), and
\begin{equation}
 \PE(\mathscr X)=\PF(\mathscr X)
 \iff
 \text{every }P_f,\ f\in F\setminus E,\ \text{commutes with every}
 \ s\in\operatorname{Bal}_E(\mathscr X),
 \label{eq:saturation}
\end{equation}
where \(\operatorname{Bal}_E\) is the balanced-kernel algebra of
\cref{def:balanced-kernels} formed with the label set \(E\).  We call \(E\)
\emph{saturated} in \(F\) when this holds.
\end{theorem}

\begin{proof}
The containment is clear.  For the criterion, \(\PF\) is generated
over \(\PE\) by the operators \(P_f\) and \(P_f^*\) with \(f\in
F\setminus E\), so equality holds if and only if every such \(P_f\) lies in
\(\PE\); the adjoints then follow because \(\PE\) is adjoint closed.  By
\cref{thm:rational-bicommutant-structural} and
\cref{thm:balanced-commutant},
\[
 \PE=(\PE)''=\operatorname{Bal}_E(\mathscr X)',
\]
and membership in the commutant of \(\operatorname{Bal}_E\) is the
stated commutation condition.
\end{proof}

\begin{corollary}
\label{cor:full-matrix}
\(\PE(\mathscr X)=\End_\Q(V_X)\) if and only if
\(\operatorname{Bal}_E(\mathscr X)=\Q\), that is, if and only if the only
balanced kernels are the scalar multiples of the identity kernel.
\end{corollary}

\begin{proof}
By \cref{thm:balanced-commutant} and the double commutant,
\(\PE=\operatorname{Bal}_E'\), and a subalgebra of \(\End_\Q(V_X)\) equals
\(\End_\Q(V_X)\) if and only if its commutant is the scalars.
\end{proof}

\begin{corollary}
\label{cor:multiplicity-free}
The defining module \(V_X\) is multiplicity free over
\(\PE(\mathscr X)\otimes\C\) if and only if \(\operatorname{Bal}_E(\mathscr X)\) is
commutative, and in that case
\(\dim_\Q Z(\PE)=\dim_\Q\operatorname{Bal}_E(\mathscr X)\).
\end{corollary}

\begin{proof}
With notation as in \eqref{eq:complex-isotypic-structural}, the commutant
after scalar extension is \(\bigoplus_j\End_\C(M_j)\), which is commutative
if and only if every \(m_j=1\).  Commutativity is preserved and reflected by
scalar extension, so it may be tested over \(\Q\).  When all \(m_j=1\) the
commutant and the center both have dimension \(t\).
\end{proof}

Throughout, ``\(V_G\) is multiplicity free'' means the condition of
\cref{cor:multiplicity-free}: the commutant \(\Pram(G)'\) is commutative.

First, \cref{thm:saturation} explains \cref{ex:all-positive-exponents}: for
\(G=C_5\) the label set \(E=\{5\}\) is not saturated in \(\{2,5\}\), because
\(\operatorname{Bal}_{\{5\}}\) contains the full matrix algebra of the
three-dimensional complementary summand, and the four-cycle \(\Psi_2\) does
not commute with all of it.  Hence
\(\Pram(C_5)\subsetneq\PAd(C_5)\), with dimensions \(5\) and \(7\).
Second, by \cref{cor:full-matrix} the equality \(\dim A=k(G)^2\) holds when
the weighted colored power graph admits no balanced kernel beyond
the scalars.  The rows of \cref{thm:comparison-panel} with
\(\dim A=k(G)^2\) therefore record rigidity of that graph.

\section{Good labels, the unit group, and rational characters}
\label{sec:goodlabels}

Throughout, \(G\) is a finite group, \(N=\exp(G)\), and
\[
 U_N=(\Z/N)^\times .
\]
We call \(n\geq1\) a \emph{good label} when \(\gcd(n,N)=1\) and a
\emph{ramified label} when every prime divisor of \(n\) divides \(N\).

\subsection{The unit group acts by isometries}

\begin{lemma}
\label{lem:good-labels}
Let \(n\) be a good label and let \(m\) satisfy \(mn\equiv1\pmod N\).  Then:
\begin{enumerate}
\item \(\Psi_n\) depends only on the residue \(n\bmod N\);
\item \(\sigma_n:\Conj(G)\to\Conj(G)\) is a bijection with inverse
\(\sigma_m\), and \(|\sigma_n(C)|=|C|\) for every class \(C\);
\item \(\Psi_n\) is orthogonal for the class form
\eqref{eq:coordinate-class-form}, and
\begin{equation}
 \Psi_n^*=\Psi_n^{-1}=\Psi_m;
 \label{eq:good-adjoint-inverse}
\end{equation}
\item the assignment \(u\mapsto\Psi_u\) is a group homomorphism
\[
 \Psi:U_N\longrightarrow O(V_G,\langle\ ,\ \rangle_G).
\]
\end{enumerate}
\end{lemma}

\begin{proof}
(1)  For \(x\in G\) the order \(|x|\) divides \(N\), so \(x^n\) depends only
on \(n\bmod|x|\), hence only on \(n\bmod N\).

(2)  Since \(x^N=1\) and \(nm\equiv1\pmod N\), we get \((x^n)^m=x^{nm}=x\)
for every \(x\in G\).  Hence \(\sigma_m\sigma_n=\sigma_n\sigma_m=\mathrm{id}\)
and \(\sigma_n\) is a bijection.  For the class sizes, fix a class \(C\) and
consider \(\theta:C\to\sigma_n(C)\), \(\theta(x)=x^n\).  It is well defined
because \(C\) generates a single class of \(n\)th powers.  It is injective:
if \(x^n=y^n\) with \(x,y\in C\) then \(x=(x^n)^m=(y^n)^m=y\).  It is
surjective: if \(z\in\sigma_n(C)\) then \(z^m\in\sigma_m\sigma_n(C)=C\) and
\(\theta(z^m)=z^{mn}=z\).  Hence \(|C|=|\sigma_n(C)|\).

(3)  By \eqref{eq:power-matrix}, \(\Psi_n\) is the permutation operator of
\(\sigma_n\).  Using (2),
\[
 \langle\Psi_nf,\Psi_nh\rangle_G
 =\frac1{|G|}\sum_{C}|C|\,f(\sigma_nC)h(\sigma_nC)
 =\frac1{|G|}\sum_{C'}|\sigma_n^{-1}C'|\,f(C')h(C')
 =\langle f,h\rangle_G,
\]
because \(|\sigma_n^{-1}C'|=|C'|\).  Thus \(\Psi_n\) is orthogonal and
\(\Psi_n^*=\Psi_n^{-1}\).  Finally \(\Psi_m\Psi_n=\Psi_{mn}=\Psi_1=I\) by (1)
and multiplicativity, which gives \eqref{eq:good-adjoint-inverse}.

(4)  Immediate from (1) and (3).
\end{proof}

\begin{proposition}
\label{prop:good-in-commutant}
For every \(u\in U_N\) we have \(\Psi_u\in\Pram(G)'\).  Consequently
\(\Psi(U_N)\) is a finite abelian group of isometries inside
\(\Pram(G)'\).
\end{proposition}

\begin{proof}
Let \(d\mid N\).  Multiplicativity gives
\(\Psi_u\Psi_d=\Psi_{ud}=\Psi_{du}=\Psi_d\Psi_u\).  Taking adjoints of this
identity and using \eqref{eq:good-adjoint-inverse},
\[
 \Psi_d^*\Psi_{u^{-1}}=\Psi_{u^{-1}}\Psi_d^*.
\]
As \(u\mapsto u^{-1}\) is a bijection of \(U_N\), the operator \(\Psi_u\)
commutes with every \(\Psi_d^*\) as well.  By
\cref{prop:minimal-prime-generators} the operators \(\Psi_d,\Psi_d^*\) with
\(d\mid N\) generate \(\Pram(G)\), and an operator commuting with a
generating set commutes with the generated algebra.
\end{proof}

\begin{proposition}
\label{prop:label-splitting}
Every \(n\geq1\) factors uniquely as \(n=n_{\mathrm{r}}n_{\mathrm{g}}\) with
\(n_{\mathrm{r}}\) a ramified label and \(n_{\mathrm{g}}\) a good label, and
then
\begin{equation}
 \Psi_n=\Psi_{n_{\mathrm{r}}}\Psi_{n_{\mathrm{g}}},
 \qquad
 \Psi_{n_{\mathrm{r}}}\in\Pram(G),
 \qquad
 \Psi_{n_{\mathrm{g}}}=\Psi_{u},\ u=n_{\mathrm{g}}\bmod N\in U_N.
 \label{eq:label-splitting}
\end{equation}
Hence
\begin{equation}
 \PAd(G)=\Alg_\Q\bigl(\Pram(G)\cup\Psi(U_N)\bigr).
 \label{eq:full-adams-generated}
\end{equation}
\end{proposition}

\begin{proof}
Let \(n_{\mathrm{r}}\) be the product of the prime powers of \(n\) at primes
dividing \(N\) and \(n_{\mathrm{g}}=n/n_{\mathrm{r}}\); uniqueness is clear.
Multiplicativity gives the first identity, and
\cref{prop:minimal-prime-generators} places \(\Psi_{n_{\mathrm{r}}}\) in
\(\Pram(G)\).  Since \(\Psi(U_N)\) is closed under adjoints by
\eqref{eq:good-adjoint-inverse}, and \(\PAd(G)\) is generated by all
\(\Psi_n\) and their adjoints, \eqref{eq:full-adams-generated} follows.
\end{proof}

\subsection{Rational classes and fields of character values}

Identify \(U_N\) with \(\operatorname{Gal}(\Q(\zeta_N)/\Q)\) by
\(u\mapsto(\zeta_N\mapsto\zeta_N^u)\).  For a rational class
\(R\in\Conj(G)/U_N\) put
\[
 \Gamma_R=\{u\in U_N:\sigma_u(C)=C\ \text{for some, equivalently every, }C\in R\},
 \qquad
 F_R=\Q(\zeta_N)^{\Gamma_R},
\]
the stabilizer being independent of the choice of \(C\) because \(U_N\) is
abelian.  The field \(F_R\) is a classical object: it is the field of
character values on \(R\).  Indeed, for \(g\in C\in R\) and
\(\chi\in\Irr(G)\), the Galois element \(u\) sends \(\chi(g)\) to
\(\chi(g^u)\), so \(u\) fixes every \(\chi(g)\) if and only if
\(\chi(g^u)=\chi(g)\) for all \(\chi\), that is, if and only if
\(g^u\) is conjugate to \(g\)~\cite{Isaacs}.  Hence
\(F_R=\Q(\chi(g):\chi\in\Irr(G))\), and the étale algebra below is the
product of the fields of values of the rational classes.  Set
\begin{equation}
 \mathcal E(G)=\prod_{R\in\Conj(G)/U_N}F_R .
 \label{eq:etale-algebra}
\end{equation}

\begin{theorem}
\label{thm:rational-class-decomposition}
There is an isomorphism of \(\Q[U_N]\)-modules
\begin{equation}
 V_G\;\cong\;\bigoplus_R\Q[U_N/\Gamma_R]\;\cong\;\bigoplus_RF_R=\mathcal E(G),
 \label{eq:etale-decomposition}
\end{equation}
under which \(\Psi_u\) acts as the field automorphism determined by \(u\) on
each factor.  In particular \([F_R:\Q]=|R|\) and
\(\dim_\Q\mathcal E(G)=k(G)\).
\end{theorem}

\begin{proof}
By \cref{lem:good-labels}, \(\Psi_u\) is the permutation operator of
\(\sigma_u\) in the basis \((\delta_C)_C\), so \(V_G\) is the permutation
module \(\Q[\Conj(G)]\) for \(U_N\).  Decomposing \(\Conj(G)\) into orbits
gives \(V_G\cong\bigoplus_R\Q[R]\), and \(\Q[R]\cong\Q[U_N/\Gamma_R]\) because
a transitive action with point stabilizer \(\Gamma_R\) is isomorphic to the
coset action.  Thus \(|R|=[U_N:\Gamma_R]=[F_R:\Q]\).

By the normal basis theorem \(\Q(\zeta_N)\cong\Q[U_N]\) as
\(\Q[U_N]\)-modules.  Taking \(\Gamma_R\)-invariants on both sides and using
that \(U_N\) is abelian, so that \(\Gamma_R\)-invariants of \(\Q[U_N]\) are
spanned by the \(\Gamma_R\)-coset sums, gives
\(F_R=\Q(\zeta_N)^{\Gamma_R}\cong\Q[U_N/\Gamma_R]\) as \(\Q[U_N]\)-modules.
Summing over \(R\) and adding the class counts gives
\(\dim_\Q\mathcal E(G)=\sum_R|R|=k(G)\).
\end{proof}

\begin{corollary}
\label{cor:good-euler-product}
For \(\operatorname{Re}s>1\),
\begin{equation}
 \prod_{p\nmid N}\det\bigl(I-\Psi_pp^{-s}\bigr)^{-1}
 =\prod_R\zeta^{(N)}_{F_R}(s)
 =\zeta^{(N)}_{\mathcal E(G)}(s),
 \label{eq:good-euler-product}
\end{equation}
the superscript denoting omission of the Euler factors at the primes dividing
\(N\).
\end{corollary}

\begin{proof}
Under \eqref{eq:etale-decomposition} the representation of
\(\operatorname{Gal}(\Q(\zeta_N)/\Q)\) on \(V_G\) is
\(\bigoplus_R\operatorname{Ind}_{\Gamma_R}^{U_N}\mathbf 1\), whose Artin
\(L\)-function is \(\prod_R\zeta_{F_R}\) by Artin induction
\cite[\S{}VII]{Washington}.  For \(p\nmid N\) the class of \(p\) in \(U_N\) is
the Frobenius at \(p\), and \(\Psi_p=\Psi_{p\bmod N}\) by
\cref{lem:good-labels}(1), so \(\det(I-\Psi_pp^{-s})^{-1}\) is the Euler
factor at \(p\) of that \(L\)-function.
\end{proof}

\begin{theorem}
\label{thm:dirichlet-factorization}
Let \(\widehat{U_N}=\operatorname{Hom}(U_N,\C^\times)\), and for
\(\chi\in\widehat{U_N}\) put
\begin{equation}
  m_G(\chi) = \#\{R\in\Conj(G)/U_N:\Gamma_R\subseteq\ker\chi\}.
  \label{eq:dirichlet-multiplicity}
\end{equation}
Then the good-label representation has the exact complex decomposition
\begin{equation}
 V_G\otimes_\Q\C
 \cong
 \bigoplus_{\chi\in\widehat{U_N}}\chi^{\oplus m_G(\chi)},
 \label{eq:dirichlet-character-decomposition}
\end{equation}
and its Euler product factors as
\begin{equation}
 \prod_{p\nmid N}\det(I-\Psi_pp^{-s})^{-1}
 =
 \prod_{\chi\in\widehat{U_N}}
 L^{(N)}(s,\chi)^{m_G(\chi)}.
 \label{eq:dirichlet-factorization}
\end{equation}
Here \(\chi\) is viewed as the corresponding Dirichlet character of
conductor dividing \(N\), and \(L^{(N)}\) denotes the Euler product with
the primes dividing \(N\) omitted.  Thus the rational-class stabilizers
determine every Dirichlet factor and its multiplicity without a
Wedderburn computation.
\end{theorem}

\begin{proof}
By \cref{thm:rational-class-decomposition}, the summand attached to \(R\)
is \(\operatorname{Ind}_{\Gamma_R}^{U_N}\mathbf1\).  Since \(U_N\) is
abelian, Frobenius reciprocity gives
\[
 \dim_\C\operatorname{Hom}_{U_N}
 \bigl(\chi,\operatorname{Ind}_{\Gamma_R}^{U_N}\mathbf1\bigr)
 =
 \dim_\C\operatorname{Hom}_{\Gamma_R}
 \bigl(\chi|_{\Gamma_R},\mathbf1\bigr)
 =
 \begin{cases}
 1,&\Gamma_R\subseteq\ker\chi,\\
 0,&\text{otherwise}.
 \end{cases}
\]
Summing over \(R\) proves
\eqref{eq:dirichlet-character-decomposition} and
\eqref{eq:dirichlet-multiplicity}.  For \(p\nmid N\), the operator
\(\Psi_p\) acts on the \(\chi\)-line by \(\chi(p)\); taking determinants
and multiplying the factors \((1-\chi(p)p^{-s})^{-1}\) over the good
primes proves \eqref{eq:dirichlet-factorization}.
\end{proof}

\begin{corollary}
\label{cor:archimedean-factor}
The field \(F_R\) is totally real if and only if \(-1\in\Gamma_R\), that is,
if and only if the classes in \(R\) are real.  Writing \((r_1(R),r_2(R))\) for
the signature of \(F_R\), the completed product
\(\prod_R\zeta_{F_R}\) has archimedean factor
\begin{equation}
 \Gamma_{\mathbf R}(s)^{\sum_R(r_1(R)+r_2(R))}\,
 \Gamma_{\mathbf R}(s+1)^{\sum_Rr_2(R)},
 \qquad
 \Gamma_{\mathbf R}(s)=\pi^{-s/2}\Gamma(s/2),
 \label{eq:archimedean-factor}
\end{equation}
and conductor \(\prod_R|d_{F_R}|\).
\end{corollary}

\begin{proof}
Complex conjugation on \(\Q(\zeta_N)\) is \(\Psi_{-1}\), and \(F_R\) is fixed
pointwise by it if and only if \(-1\in\Gamma_R\); a subfield of a cyclotomic
field is either totally real or totally imaginary, so this is the criterion
for total reality.  By \cref{lem:good-labels}(2) the condition
\(\sigma_{-1}(C)=C\) says that \(C=C^{-1}\).  The displayed factor is
\(\prod_R\Gamma_{\mathbf R}(s)^{r_1(R)}\Gamma_{\mathbf C}(s)^{r_2(R)}\) with
\(\Gamma_{\mathbf C}(s)=\Gamma_{\mathbf R}(s)\Gamma_{\mathbf R}(s+1)\).
\end{proof}

For the Monster, \(k(\M)=194\) with \(172\) rational classes
(\cref{prop:monster-imaginary}), so \(150\) are singletons and \(22\) have two
elements, the latter imaginary quadratic by
\eqref{eq:monster-discriminants}.  Then \(\sum_R(r_1+r_2)=172\),
\(\sum_Rr_2=22\), and \eqref{eq:archimedean-factor} reads
\[
 \Gamma_{\mathbf R}(s)^{172}\,\Gamma_{\mathbf R}(s+1)^{22}.
\]
Here every \(F_R\) has degree at most two.  For \(J_1\) the orbit of size
three at element order \(19\) has \(\Gamma_R\) of order six in
\(U_{19}\cong\Z/18\), and that subgroup contains the involution, so \(F_R\) is
a totally real cubic field contributing \(\Gamma_{\mathbf R}(s)^3\).  Its
block center \(\Q(\zeta_3)\) in \cref{thm:regular-rational-class} is another
field: \(\Q[U_N/\Gamma_R]\cong\Q\oplus W_R\) with \(F_R\) on the left and
\(\Q(\zeta_3)=\End_{\Q[U_N]}(W_R)\) on the right.

\subsection{Ramified labels}

\begin{theorem}
\label{thm:ramified-spectrum}
Let \(p\mid N\), let \(b=v_p(N)\), and let \(\Conj(G)_{p'}\) be the set of
classes of elements of order prime to \(p\).  Then \(\sigma_p\) restricts to a
bijection of \(\Conj(G)_{p'}\), the span
\[
 W_p=\Span_\Q\bigl(\delta_C:C\notin\Conj(G)_{p'}\bigr)
\]
is \(\Psi_p\)-invariant with \(\Psi_p^{\,b}|_{W_p}=0\), and \(\Psi_p\) induces
on \(V_G/W_p\) the permutation operator of \(\sigma_p|_{\Conj(G)_{p'}}\).
Consequently the characteristic polynomial of \(\Psi_p\) is
\[
 X^{\dim W_p}\prod_{\gamma}\bigl(X^{\ell(\gamma)}-1\bigr),
\]
the product being over the cycles \(\gamma\) of
\(\sigma_p|_{\Conj(G)_{p'}}\), and every nonzero eigenvalue of \(\Psi_p\) is a
root of unity.
\end{theorem}

\begin{proof}
If \(|x|=m\) with \(p\nmid m\) then \(|x^p|=m\), so \(\sigma_p\) maps
\(\Conj(G)_{p'}\) into itself; it is injective there, since choosing \(m'\)
with \(pm'\equiv1\bmod m\) recovers \(x\) from \(x^p\) inside
\(\langle x\rangle\), hence bijective.  Conversely if \(p\mid|x|\) then
\(p\mid|x^p|\) unless \(|x|=p\), in which case \(x^p=1\); so for
\(C\notin\Conj(G)_{p'}\) every class in \(\sigma_p^{-1}(C)\) also lies outside
\(\Conj(G)_{p'}\).  By \eqref{eq:power-matrix},
\(\Psi_p\delta_C=\sum_{\sigma_p(C')=C}\delta_{C'}\), so \(W_p\) is
\(\Psi_p\)-invariant.  Moreover \(|x^{p^{b}}|\) is prime to \(p\) for every
\(x\in G\), so \(\sigma_p^{\,b}(\Conj(G))\subseteq\Conj(G)_{p'}\) and
\(\sigma_p^{-b}(C)=\emptyset\) for \(C\notin\Conj(G)_{p'}\); hence
\(\Psi_p^{\,b}|_{W_p}=0\).

The images of the \(\delta_C\) with \(C\in\Conj(G)_{p'}\) form a basis of
\(V_G/W_p\), and \(\Psi_p\) sends \(\delta_C\) to
\(\delta_{\sigma_p^{-1}(C)\cap\Conj(G)_{p'}}\) modulo \(W_p\), which is the
permutation operator of the bijection of the first paragraph.  The
characteristic polynomial is the product of those of \(\Psi_p|_{W_p}\) and of
the induced map, and a permutation of cycle type \((\ell(\gamma))_\gamma\) has
characteristic polynomial \(\prod_\gamma(X^{\ell(\gamma)}-1)\).
\end{proof}

\subsection{Local factors at ramified primes}
\label{sec:ramified-completion}

The good primes give the Euler factors of \(\zeta_{\mathcal E(G)}\) away
from \(N\).  The factors at \(p\mid N\) are determined, as for any
abelian étale algebra, by Frobenius acting on inertia invariants; here both
are built from good-label operators restricted to element-order layers.

For \(m\mid N\), let \(X_m\) be the classes of elements of order \(m\)
and put \(V_m=\operatorname{Map}(X_m,\Q)\).  The action of \(U_N\) on
\(X_m\) factors through \(U_m=(\Z/m)^\times\).  We write
\(\Psi_{u,m}\) for the resulting operator; a lift of \(u\) to \(U_N\)
gives the same operator on \(V_m\).

Fix a prime \(p\).  If \(p\nmid m\), set \(I_{p,m}=1\) and
\(\kappa_{p,m}=p\bmod m\).  If \(m=p^am'\), with \(a>0\) and
\(p\nmid m'\), put
\[
 I_{p,m}=\ker(U_m\longrightarrow U_{m'})
\]
and choose \(\kappa_{p,m}\in U_m\) such that
\begin{equation}
 \kappa_{p,m}\equiv p\pmod{m'},\qquad
 \kappa_{p,m}\equiv1\pmod{p^a}.
 \label{eq:k-p-conditions}
\end{equation}
Let
\begin{equation}
 \Pi_{p,m}=\frac1{|I_{p,m}|}
 \sum_{u\in I_{p,m}}\Psi_{u,m},
 \qquad
 \Phi_{p,m}
 =\Pi_{p,m}\Psi_{\kappa_{p,m},m}\Pi_{p,m}
 \bigm|_{V_m^{I_{p,m}}}.
 \label{eq:frobenius-operator}
\end{equation}
The first operator is the orthogonal projection onto
\(V_m^{I_{p,m}}\).  The second is independent of the choice in
\eqref{eq:k-p-conditions}, since two choices differ by an element of
\(I_{p,m}\).

Define
\begin{equation}
 \mathcal L_p(G,T)=
 \prod_{m\mid N}
 \det\bigl(I-\Phi_{p,m}T\mid V_m^{I_{p,m}}\bigr).
 \label{eq:complete-local-polynomial}
\end{equation}

\begin{theorem}
\label{thm:ramified-completion}
For \(\operatorname{Re}s>1\),
\begin{equation}
 \zeta_{\mathcal E(G)}(s)
 =\prod_R\zeta_{F_R}(s)
 =\prod_p\mathcal L_p(G,p^{-s})^{-1}.
 \label{eq:complete-euler-product}
\end{equation}
If \(p\nmid N\), then
\[
 \mathcal L_p(G,T)=\det(I-\Psi_pT\mid V_G).
\]
If \(p\mid N\), the factors with \(p\nmid m\) are the Frobenius
determinants on the \(p\)-regular order layers, and
\begin{equation}
 C_p(G,T)=
 \prod_{\substack{m\mid N\\p\mid m}}
 \det\bigl(I-\Phi_{p,m}T\mid V_m^{I_{p,m}}\bigr)
 \label{eq:Cp-polynomial}
\end{equation}
is the contribution of the \(p\)-singular layers.
\end{theorem}

\begin{proof}
Every \(U_N\)-orbit \(R\) lies in one \(X_m\).  On
\(\Q[R]\), the local Artin factor of \(F_R\) at \(p\) is
\[
 \det\bigl(I-\operatorname{Frob}_pT
 \mid\Q[R]^{I_{p,m}}\bigr)^{-1}.
\]
For \(p\nmid m\), inertia is trivial and \(p\bmod m\) acts as Frobenius.
For \(p\mid m\), \eqref{eq:k-p-conditions} represents Frobenius modulo
inertia, so \(\Phi_{p,m}\) is its action on the inertia invariants.
Multiplication over the orbits in \(X_m\), then over \(m\mid N\), gives
\(\mathcal L_p(G,T)^{-1}\).  Multiplication over \(p\) gives
\eqref{eq:complete-euler-product}.  If \(p\nmid N\), all inertia groups
are trivial and the direct sum of the \(\Psi_{p,m}\) is \(\Psi_p\).
\end{proof}

\begin{corollary}
\label{cor:complete-euler-product}
The class sizes together with all power maps, that is, the ramified eidetic
object \(\EramObj(G)\), determine every local factor
of \(\zeta_{\mathcal E(G)}\).
\end{corollary}

\begin{corollary}
\label{cor:ramified-degree}
The degree of \(C_p(G,T)\) is
\[
 \delta_p(G)=
 \sum_{\substack{m\mid N\\p\mid m}}
 \#(I_{p,m}\backslash X_m).
\]
If \(r_p(G)\) is the number of \(p\)-singular rational classes and
\(c_p(G)\) the number of \(p\)-singular conjugacy classes, then
\[
 r_p(G)\leq\delta_p(G)\leq c_p(G).
\]
\end{corollary}

\begin{proof}
The fixed space of a permutation group on \(X_m\) has dimension equal to
its number of orbits.  Since \(I_{p,m}\leq U_m\), its orbit partition
refines the rational-class partition and coarsens the partition into
points.
\end{proof}

For the Monster, the ramified degrees are
\[
\begin{array}{c|rrrrrrrrrrrrrrr}
p&2&3&5&7&11&13&17&19&23&29&31&41&47&59&71\\
\hline
\delta_p&
130&90&44&28&11&12&5&4&5&2&3&1&2&1&1
\end{array}
\]
At \(p=2,3\),
\[
 122<130<133,\qquad 88<90<93.
\]

\subsection{Good-label representations}

\begin{theorem}
\label{thm:multiplicity-good-labels}
Let \(A=\Pram(G)\) and \(N=\exp(G)\).  Choose representatives
\(S_1,\ldots,S_t\) for the simple rational \(A\)-modules occurring in
\(V_G\), and write
\[
 D_j=\End_A(S_j),\qquad
 M_j=\operatorname{Hom}_A(S_j,V_G).
\]
Then evaluation gives an \(A\)-module decomposition
\begin{equation}
 V_G\cong\bigoplus_{j=1}^t S_j\otimes_{D_j}M_j,
 \label{eq:double-centralizer-decomposition}
\end{equation}
and restriction to the isotypic components gives
\begin{equation}
 A'\cong\bigoplus_{j=1}^t\End_{D_j}(M_j).
 \label{eq:double-centralizer-commutant}
\end{equation}
For every \(j\), the good labels induce a canonical representation
\begin{equation}
 \rho_j:U_N\longrightarrow
 \End_{D_j}(M_j)^\times,\qquad
 \rho_j(u)(f)=\Psi_u\circ f ,
 \label{eq:good-label-representation}
\end{equation}
with the following properties.
\begin{enumerate}
\item The image is finite and
  \(\rho_j(u)^\dagger=\rho_j(u^{-1})\), where \(\dagger\) is the positive
  involution on \(\End_{D_j}(M_j)\) induced by the class form.
\item Every complex eigenvalue of every \(\rho_j(u)\) is a root of unity,
  and every rational characteristic polynomial of \(\rho_j(u)\) is a
  product of cyclotomic polynomials.
\item The action on the \(j\)-th component is scalar if and only if
  \(\rho_j(U_N)\subseteq Z\bigl(\End_{D_j}(M_j)\bigr)^\times\).
  In that event it defines a character
  \(\lambda_j:U_N\to\mu(K_j)\), where \(K_j=Z(D_j)\).
\item If \(D_j=K_j\) and \(\dim_{K_j}M_j=1\), the action is automatically
  scalar.  Thus the split multiplicity-one theorem below is the rank-one
  specialization of \eqref{eq:good-label-representation}.
\item Put \(a_j=\dim_{D_j}S_j\).  For every \(u\in U_N\),
  \begin{equation}
  \#\{C\in\Conj(G):\sigma_u(C)=C\}
  = \sum_{j=1}^t a_j \operatorname{Tr}_\Q(\rho_j(u)\mid M_j).
  \label{eq:general-lefschetz}
  \end{equation}
\end{enumerate}
\end{theorem}

\begin{proof}
The algebra \(A\) is semisimple by
\cref{thm:rational-bicommutant-structural}.  The rational
double-centralizer theorem therefore gives
\eqref{eq:double-centralizer-decomposition} and
\eqref{eq:double-centralizer-commutant}.  Explicitly, \(S_j\) is a right
\(D_j\)-module through evaluation by \(A\)-endomorphisms, \(M_j\) is the
corresponding left \(D_j\)-multiplicity module, and the evaluation map
\[
 S_j\otimes_{D_j}\operatorname{Hom}_A(S_j,V_G)\longrightarrow V_G,\qquad
 s\otimes f\longmapsto f(s),
\]
identifies the source with the \(S_j\)-isotypic component.

Every \(\Psi_u\) commutes with \(A\) by \cref{prop:good-in-commutant}, so
postcomposition preserves \(\operatorname{Hom}_A(S_j,V_G)\) and is
\(D_j\)-linear.  The equality \(\Psi_{uv}=\Psi_u\Psi_v\) proves that
\(\rho_j\) is a representation.  Because \(U_N\) is finite, its image is
finite.  The identity \(\Psi_u^*=\Psi_{u^{-1}}\) of
\eqref{eq:good-adjoint-inverse} restricts under the double-centralizer
identification to \(\rho_j(u)^\dagger=\rho_j(u^{-1})\), proving (1).

If \(u\) has order \(e\) in \(U_N\), then \(\rho_j(u)^e=1\).  In
characteristic zero the minimal polynomial divides \(x^e-1\), which is
square-free, so \(\rho_j(u)\) is semisimple with roots of unity as
eigenvalues.  Its rational characteristic polynomial is therefore a product
of cyclotomic polynomials.  This proves (2).

For (3),
\[
 Z\bigl(\End_{D_j}(M_j)\bigr)=Z(D_j)=K_j,
\]
and every finite subgroup of \(K_j^\times\) lies in \(\mu(K_j)\).  If
\(D_j=K_j\) and \(\dim_{K_j}M_j=1\), then
\(\End_{K_j}(M_j)=K_j\), proving (4).

Finally, \(\Psi_u\) is the permutation operator of the \(u\)-power action
on conjugacy classes, so its rational trace is the number of fixed classes.
Choose a right \(D_j\)-basis of \(S_j\) of size \(a_j\).  On
\(S_j\otimes_{D_j}M_j\), the operator induced by \(\Psi_u\) is the direct
sum of \(a_j\) copies of \(\rho_j(u)\) as a rational linear operator.
Taking traces in \eqref{eq:double-centralizer-decomposition} gives
\eqref{eq:general-lefschetz}.
\end{proof}

For \(m_j>1\), the good labels define a finite \(\dagger\)-unitary
representation in \(\End_{D_j}(M_j)\).  A scalar character exists if its
image lies in the center.

\subsection{Rational characters on the isotypic components}

Write \(A=\Pram(G)\).  \Cref{prop:good-in-commutant} says that the good
labels live in \(A'\); the structure theory of \cref{sec:structural-theory}
therefore constrains them as soon as we know \(A'\).  The strongest
statement occurs in the multiplicity-free case, which by
\cref{cor:multiplicity-free} is the case \(A'\) commutative.

\begin{theorem}
\label{thm:good-label-characters}
Suppose \(V_G\) is multiplicity free over \(A=\Pram(G)\).  Let
\(E_1,\dots,E_t\) be the primitive idempotents of \(Z(A)\), let
\(A E_k\cong M_{r_k}(D_k)\) be the corresponding simple factors, and let
\(K_k=Z(AE_k)\) be their center fields.  Then \(A'=Z(A)\), and for every
\(u\in U_N\)
\begin{equation}
 \Psi_uE_k=\lambda_k(u)\,E_k,
 \qquad
 \lambda_k:U_N\longrightarrow\mu(K_k),
 \label{eq:good-label-character}
\end{equation}
where \(\mu(K_k)\) is the group of roots of unity of \(K_k\) and
\(\lambda_k\) is a group homomorphism.  Its conductor divides \(N\).  In
particular:
\begin{enumerate}
\item if \(K_k=\Q\) then \(\lambda_k\) is either trivial or a quadratic
Dirichlet character modulo \(N\);
\item for every prime \(p\nmid N\), the compression
\(A_k(p)=\Psi_p|_{V_k}\) satisfies
\begin{equation}
 A_k(p)=\lambda_k(p)\,I_{V_k},
 \qquad
 \det\bigl(I-A_k(p)T\bigr)=\prod_{\iota}\bigl(1-\iota(\lambda_k(p))T\bigr)^{d_k/[K_k:\Q]},
 \label{eq:good-local-factor}
\end{equation}
the product running over the embeddings \(\iota:K_k\hookrightarrow\C\) and
\(d_k=\dim_\Q V_k\).  When \(K_k=\Q\) this reduces to
\(\det(I-A_k(p)T)=(1-\lambda_k(p)T)^{d_k}\).
\end{enumerate}
\end{theorem}

\begin{proof}
By \cref{cor:multiplicity-free} the commutant \(A'\) is commutative and
\(\dim_\Q A'=\dim_\Q Z(A)\).  Since \(Z(A)\subseteq A'\) always, the two
coincide, so \(A'=Z(A)\subseteq A\).

Fix \(u\in U_N\).  By \cref{prop:good-in-commutant}, \(\Psi_u\in A'=Z(A)\).
Multiplying by the primitive central idempotent \(E_k\) gives an element of
\(Z(AE_k)=K_k\), which is \eqref{eq:good-label-character} once we note that
\(K_k\) acts on the ideal \(AE_k\) through the scalar \(\lambda_k(u)\).  The
map \(\lambda_k\) is multiplicative because \(u\mapsto\Psi_u\) is
(\cref{lem:good-labels}(4)) and \(E_k\) is idempotent.  Since \(U_N\) is
finite, \(\lambda_k(u)\) is a root of unity in \(K_k\).  The conductor
divides \(N\) because \(\lambda_k\) is defined on \(U_N\) and
\(\Psi_n\) depends only on \(n\bmod N\) by \cref{lem:good-labels}(1).

(1)  The only roots of unity in \(\Q\) are \(\pm1\), so a homomorphism
\(U_N\to\mu(\Q)=\{\pm1\}\) is trivial or of order two; a homomorphism
\(U_N\to\{\pm1\}\) is a real Dirichlet character modulo \(N\), and
a nontrivial one is quadratic.

(2)  By \eqref{eq:good-label-character}, \(\Psi_p\) acts on
\(V_k=E_kV_G\) as multiplication by the scalar \(\lambda_k(p)\in K_k\).  The
characteristic polynomial of multiplication by \(\alpha\in K_k\) on a
\(K_k\)-vector space of \(\Q\)-dimension \(d_k\) is
\(\prod_\iota(T-\iota(\alpha))^{d_k/[K_k:\Q]}\), which gives
\eqref{eq:good-local-factor}.
\end{proof}

\subsection{Split multiplicity-one blocks}

\begin{theorem}
\label{thm:block-local-character}
Let \(A=\Pram(G)\), let \(E\) be a primitive central idempotent of \(A\) with
center field \(K=Z(AE)\).  Assume that the simple factor is split over its
center, \(AE\cong M_r(K)\), and let \(m\) be the multiplicity of its natural
simple module in \(V_G\).  If \(m=1\), then for every \(u\in U_N\)
\begin{equation}
 \Psi_uE=\lambda_E(u)\,E,
 \qquad
 \lambda_E:U_N\longrightarrow\mu(K),
 \label{eq:block-local-character}
\end{equation}
and \(\lambda_E\) is a homomorphism of conductor dividing \(N\).  No
hypothesis is placed on the other blocks.
\end{theorem}

\begin{proof}
This is the rank-one specialization of
\cref{thm:multiplicity-good-labels}.  Splitness \(AE\cong M_r(K)\) gives
\(D_j=K_j=K\) on the corresponding component, and multiplicity one gives
\(\dim_{K}M_j=1\); by parts (3) and (4) the representation \(\rho_j\) is a
scalar character \(\lambda_E:U_N\to\mu(K)\), which is
\eqref{eq:block-local-character}.  The conductor divides \(N\) because
\(\Psi_n\) depends only on \(n\bmod N\) by \cref{lem:good-labels}(1).
\end{proof}

For a general simple factor \(AE\cong M_r(D)\) with center \(K\), the
double-centralizer theorem gives \(EA'E\cong M_m(D^{\mathrm{op}})\).
For \(m=1\), the good-label action lies in \(D^{\mathrm{op}}\).  If
\(D=K\), it is scalar.

\subsection{The involution on the center}

\begin{lemma}
\label{lem:star-fixes-idempotents}
Every primitive central idempotent \(E\) of \(A=\Pram(G)\) satisfies
\(E^*=E\), and the star restricts to an involution
\(\iota_E\) of the field \(K=Z(AE)\).
\end{lemma}

\begin{proof}
The class form \eqref{eq:coordinate-class-form} is positive definite and
\(A\) is closed under its adjoint, so the isotypic decomposition of \(V_G\)
under \(A\) is orthogonal and \(E\) is the orthogonal projection onto its
isotypic component; an orthogonal projection is self-adjoint.  The star is an
anti-automorphism of \(A\) fixing \(E\), so it preserves the ideal \(AE\) and
hence its center \(K\), on which an anti-automorphism is an automorphism;
it is an involution because the star is.
\end{proof}

\begin{proposition}
\label{prop:star-inverts}
With the hypotheses of \cref{thm:block-local-character},
\begin{equation}
 \lambda_E(u^{-1})=\iota_E\bigl(\lambda_E(u)\bigr)
 =\lambda_E(u)^{-1}
 \qquad(u\in U_N).
 \label{eq:star-inverts}
\end{equation}
Consequently \(\iota_E\) acts on the image group \(\lambda_E(U_N)\) as
inversion.
\end{proposition}

\begin{proof}
Apply the star to \eqref{eq:block-local-character}.  On the left,
\((\Psi_uE)^*=E^*\Psi_u^*=E\Psi_{u^{-1}}=\lambda_E(u^{-1})E\) by
\cref{lem:star-fixes-idempotents} and \eqref{eq:good-adjoint-inverse}.  On
the right, \(\lambda_E(u)\) lies in \(K\subseteq Z(A)\), so
\((\lambda_E(u)E)^*=\iota_E(\lambda_E(u))E\).  Comparing gives the first
equality in \eqref{eq:star-inverts}; the second is multiplicativity of
\(\lambda_E\).
\end{proof}

\begin{corollary}
\label{cor:quadratic-sharp}
Let \(L_E=\Q\bigl(\lambda_E(U_N)\bigr)\subseteq K\) be the field generated by
the realized values.  The following alternatives are mutually exclusive
and exhaustive.
\begin{enumerate}
\item \(\iota_E\) is trivial on \(L_E\).  Then \(\lambda_E^2=\mathbf1\), so
\(\lambda_E\) is trivial or a quadratic Dirichlet character modulo \(N\),
and \(L_E=\Q\).
\item \(\iota_E\) is nontrivial on \(L_E\).  Then \(\lambda_E\) has order
\(e>2\), \(L_E=\Q(\zeta_e)\) is a CM field, and \(\iota_E\) restricts to its
complex conjugation.
\end{enumerate}
In particular \(K=\Q\) implies (1), but (1) is strictly weaker than
\(K=\Q\): what forces quadratic values is triviality of the star on the
generated value field, not rationality of the center.
\end{corollary}

\begin{proof}
By \eqref{eq:star-inverts}, \(\iota_E\) inverts every value.  If \(\iota_E\)
is trivial on \(L_E\) then \(\lambda_E(u)=\lambda_E(u)^{-1}\) for all \(u\),
so every value is \(\pm1\) and \(L_E=\Q\); \cref{thm:good-label-characters}(1)
then identifies \(\lambda_E\).  Otherwise the image is cyclic of some order
\(e\), generated by a primitive \(e\)th root of unity \(\zeta\) with
\(L_E=\Q(\zeta_e)\), and \(\iota_E(\zeta)=\zeta^{-1}\) is a nontrivial
automorphism, so \(e>2\).  A nontrivial automorphism of \(\Q(\zeta_e)\)
inverting \(\zeta_e\) is complex conjugation in every embedding, and a
cyclotomic field of degree greater than one on which complex conjugation is
nontrivial is CM.
\end{proof}

A block carries a character of order greater than two only if its center
admits a complex conjugation, and the values then generate a cyclotomic
subfield of that center.  The center fields occurring in
\cref{sec:characters} are accordingly \(\Q\), \(\Q(\zeta_3)=\Q(\sqrt{-3})\),
and \(\Q(\zeta_5)\); no real quadratic or real quartic center carries a
nontrivial character.

\subsection{Support and character order}

\begin{lemma}
\label{lem:orbit-exponent-bound}
Let \(E\) be a central idempotent of \(A=\Pram(G)\) with \(V_E=EV_G\neq0\)
on which the good labels act by scalars,
\(\Psi_uE=\lambda_E(u)E\).  Let
\[
 S(E)=\{C\in\Conj(G):f(C)\neq0\ \text{for some}\ f\in V_E\}
\]
be the class support of the block.  Then \(S(E)\) is stable under every
\(\sigma_u\), and for each \(u\in U_N\) the multiplicative order of
\(\lambda_E(u)\) divides the order of the permutation
\(\sigma_u|_{S(E)}\).  Consequently the order of \(\lambda_E\) divides the
exponent of the image of \(U_N\) in the symmetric group of \(S(E)\).  In
particular, if every \(\Psi(U_N)\)-orbit meeting \(S(E)\) has at most two
classes, then \(\lambda_E\) is trivial or quadratic.
\end{lemma}

\begin{proof}
For \(f\in V_E\) we have \(\Psi_uf=\lambda_E(u)f\) up to the idempotent:
\(\Psi_u(Ef)=\lambda_E(u)Ef\), and \(\lambda_E(u)\neq0\), so
\(\Psi_uV_E=V_E\).  Since \((\Psi_uf)(C)=f(\sigma_uC)\), the support
satisfies \(\sigma_u^{-1}S(E)\subseteq S(E)\), hence equality by
finiteness.  Now let \(m\) be the order of \(\sigma_u|_{S(E)}\) and take
\(f\in V_E\).  For \(C\in S(E)\),
\((\Psi_u^mf)(C)=f(\sigma_u^mC)=f(C)\); for \(C\notin S(E)\), both sides
vanish because \(S(E)\) is \(\sigma_u\)-stable and \(f\) is supported on
\(S(E)\).  Hence \(\Psi_u^mf=f\), so \(\lambda_E(u)^m=1\).  The exponent
statement follows by taking least common multiples over \(u\), and the
final claim holds because a transitive abelian action is regular: on an
orbit of at most two classes every \(\sigma_u\) squares to the identity,
so \(\lambda_E(u)^2=1\) for every \(u\).
\end{proof}

A character of order greater than two therefore requires a
\(\Psi(U_N)\)-orbit of at least three classes inside the block support.

\begin{proposition}
\label{prop:classical-series}
Let \(n\geq1\).
\begin{enumerate}
\item Every conjugacy class of \(S_n\) is rational.  Hence \(\Psi_u=I\) for
every \(u\in U_N\), every representation \(\rho_j\) of
\cref{thm:multiplicity-good-labels} is trivial, and every good-label
character of \(\Pram(S_n)\) is trivial.
\item Every \(\Psi(U_N)\)-orbit on \(\Conj(A_n)\) has at most two classes.
Hence every \(\rho_j\) has an elementary abelian \(2\)-group as image, and
every scalar block character of \(\Pram(A_n)\) is trivial or quadratic.
\end{enumerate}
\end{proposition}

\begin{proof}
(1)  A permutation and its coprime powers have the same cycle type, hence
are conjugate in \(S_n\).  So \(\sigma_u=\mathrm{id}\) for every good label
\(u\), and \(\Psi_u=I\).

(2)  Let \(x\in A_n\) and let \(u\) be a good label.  Then \(x^u\) has the
same cycle type as \(x\), so the \(A_n\)-class of \(x^u\) lies inside the
\(S_n\)-class of \(x\).  An \(S_n\)-class meets \(A_n\) in at most two
\(A_n\)-classes, so the \(\Psi(U_N)\)-orbit of \([x]\) has at most two
elements.  Every \(\sigma_u\) therefore has order at most two on
\(\Conj(A_n)\), so \(\Psi_u^2=I\), which forces \(\rho_j(u)^2=1\) on every
multiplicity space; the scalar statement is
\cref{lem:orbit-exponent-bound}.
\end{proof}

\subsection{A trace identity}

The following is the all-scalar case of \eqref{eq:general-lefschetz}.

\begin{proposition}
\label{prop:lefschetz}
Let \(E_1,\dots,E_t\) be the primitive central idempotents of
\(A=\Pram(G)\), with center fields \(K_k\) and \(V_k=E_kV_G\), and suppose
each acts by a center scalar, \(\Psi_uE_k=\lambda_k(u)E_k\), for instance
because every simple factor is split over its center and every block has
multiplicity one (\cref{thm:block-local-character}).  Then for every
\(u\in U_N\)
\begin{equation}
 \#\{C\in\Conj(G):\sigma_u(C)=C\}
 =\sum_{k=1}^{t}\frac{d_k}{[K_k:\Q]}\;
 \operatorname{Tr}_{K_k/\Q}\bigl(\lambda_k(u)\bigr),
 \qquad d_k=\dim_\Q V_k .
 \label{eq:lefschetz}
\end{equation}
\end{proposition}

\begin{proof}
By \eqref{eq:power-matrix} the operator \(\Psi_u\) is the permutation matrix
of \(\sigma_u\) in the coordinate basis \((\delta_C)_C\), so its trace is the
number of fixed classes.  Trace is additive over
\(V_G=\bigoplus_kV_k\), and by hypothesis the
restriction \(\Psi_u|_{V_k}\) is multiplication by \(\lambda_k(u)\in K_k\) on
a \(K_k\)-vector space of \(\Q\)-dimension \(d_k\).  The \(\Q\)-trace of
multiplication by \(\alpha\in K_k\) on such a space is
\(\bigl(d_k/[K_k:\Q]\bigr)\operatorname{Tr}_{K_k/\Q}(\alpha)\).
\end{proof}

For a rational center, \cref{thm:good-label-characters}(1) gives a trivial
or quadratic character.  Under the isomorphism
\(U_N\cong\operatorname{Gal}(\Q(\zeta_N)/\Q)\), a quadratic character
\(\lambda_k\) corresponds to a quadratic subfield
\(\Q(\sqrt{D_k})\subseteq\Q(\zeta_N)\), and \(\lambda_k\) is the Kronecker
symbol \(\bigl(\tfrac{D_k}{\cdot}\bigr)\) of the associated fundamental
discriminant \(D_k\), whose prime divisors divide \(N\)~\cite{Davenport}.

\subsection{Saturation for finite groups}

\begin{theorem}
\label{thm:good-label-saturation}
If \(V_G\) is multiplicity free over \(\Pram(G)\), then
\begin{equation}
 \Pram(G)=\PAd(G).
 \label{eq:ram-equals-ad}
\end{equation}
\end{theorem}

\begin{proof}
By \cref{thm:good-label-characters} we have \(A'=Z(A)\) for
\(A=\Pram(G)\), and \(\Psi_u\in A'=Z(A)\subseteq A\) for every \(u\in U_N\).
Hence \(\Psi(U_N)\subseteq\Pram(G)\), and
\eqref{eq:full-adams-generated} gives
\(\PAd(G)\subseteq\Pram(G)\).  The reverse inclusion is trivial.
\end{proof}

\subsection{The rational-class splitting}

Write
\(\Gamma=\Psi(U_N)\subseteq O(V_G)\) and
\[
 V_G^{\Gamma}=\{f\in V_G:\Psi_uf=f\ \text{for all}\ u\in U_N\}.
\]

\begin{proposition}
\label{prop:rational-splitting}
The subspace \(V_G^{\Gamma}\) and its class-orthogonal complement are both
reducing for \(\PAd(G)\), and
\begin{equation}
 \dim_\Q V_G^{\Gamma}=\#\{\text{rational classes of }G\}.
 \label{eq:rational-class-dimension}
\end{equation}
If moreover \(V_G\) is multiplicity free over \(A=\Pram(G)\), then
\begin{equation}
 V_G^{\Gamma}=\bigoplus_{\lambda_k=\mathbf1}V_k,
 \qquad\text{so}\qquad
 \sum_{\lambda_k=\mathbf1}d_k=\#\{\text{rational classes}\},
 \quad
 \sum_{\lambda_k\neq\mathbf1}d_k=k(G)-\#\{\text{rational classes}\}.
 \label{eq:trivial-block-sum}
\end{equation}
\end{proposition}

\begin{proof}
Every element of \(\Gamma\) commutes with \(\Pram(G)\) by
\cref{prop:good-in-commutant} and with \(\Gamma\) itself, so
\(V_G^{\Gamma}\) is invariant under all of \(\Pram(G)\) and all of
\(\Gamma\); by \eqref{eq:full-adams-generated} it is invariant under
\(\PAd(G)\).  Since \(\Gamma\) consists of isometries, the orthogonal
complement of \(V_G^\Gamma\) is \(\Gamma\)-invariant, and it is
\(\PAd(G)\)-invariant because \(\PAd(G)\) is adjoint closed and preserves
\(V_G^\Gamma\).

For \eqref{eq:rational-class-dimension}: \(\Gamma\) acts on \(V_G\) by
permuting the coordinate basis \((\delta_C)_C\) according to the
\(\sigma_u\), so the fixed space has as basis the orbit sums, one for each
\(\Gamma\)-orbit on \(\Conj(G)\).  By \cref{lem:good-labels}(2) and the
remark following it, those orbits are the rational classes.

For \eqref{eq:trivial-block-sum}: assume multiplicity freeness and fix
\(u\in U_N\).  By \eqref{eq:good-label-character}, \(\Psi_u\) acts on
\(V_k\) as multiplication by \(\lambda_k(u)\in K_k\).  If
\(\lambda_k=\mathbf1\) then \(V_k\subseteq V_G^\Gamma\).  If
\(\lambda_k\neq\mathbf1\), choose \(u\) with \(\lambda_k(u)\neq1\); then
\(\lambda_k(u)-1\) is a nonzero element of the field \(K_k\), hence
invertible, so \(\Psi_u-I\) is invertible on the \(K_k\)-module \(V_k\) and
\(V_k\cap V_G^\Gamma=0\).  As \(V_G=\bigoplus_kV_k\) and each \(V_k\) is
\(\Gamma\)-invariant, the fixed space is the sum of the trivial blocks.  The
two displayed sums then follow from \(\sum_kd_k=k(G)\).
\end{proof}

\Cref{prop:rational-splitting} is the general form of the coarse
decomposition that any concrete census exhibits: the total class count
splits as (rational classes) plus (nontrivially twisted dimension), and in
the multiplicity-free case the splitting is by blocks.  For the Monster,
where \(k(G)=194\) and there are \(172\) rational classes, it forces
\(172=\sum_{\lambda_k=\mathbf1}d_k\) and
\(22=\sum_{\lambda_k\neq\mathbf1}d_k\), and \cref{prop:monster-imaginary}
identifies the blocks on either side.

\begin{proposition}
\label{prop:real-rational-parity}
Assume that \(V_G\) is multiplicity free over \(\Pram(G)\).  Then every
nontrivial good-label character \(\lambda_k\) is odd if and only if the
number of real conjugacy classes of \(G\) equals the number of rational
classes of \(G\).  Namely,
\begin{equation}
 \#\{\text{real classes}\}-\#\{\text{rational classes}\}
 =
 \sum_{\substack{\lambda_k\neq\mathbf1\\ \lambda_k(-1)=1}}d_k .
 \label{eq:real-rational-parity}
\end{equation}
\end{proposition}

\begin{proof}
The real classes are the fixed points of inversion on
\(\Conj(G)\), so their number is the dimension of the \(+1\)-eigenspace of
\(\Psi_{-1}\).  By multiplicity freeness, \(\Psi_{-1}\) acts on \(V_k\) as
the scalar \(\lambda_k(-1)\), and hence
\[
 \#\{\text{real classes}\}
 =\sum_{\lambda_k(-1)=1}d_k .
\]
By \eqref{eq:trivial-block-sum}, the sum over the trivial characters is the
number of rational classes.  Subtracting gives
\eqref{eq:real-rational-parity}.  The right-hand side vanishes if and only
if every nontrivial character is odd.
\end{proof}

\subsection{Compressions}

Channel compressions obey the multiplicative law needed to define local
block data without an ordering convention.

\begin{proposition}
\label{prop:compression-multiplicative}
Let \(E\) be a central idempotent of \(A=\Pram(G)\), let \(V_E=EV_G\), and
for \(n\geq1\) put \(A_E(n)=E\Psi_nE|_{V_E}\).  Then \(E\) commutes with
\(\Psi_n\) for every \(n\geq1\), and
\begin{equation}
 A_E(m)A_E(n)=A_E(mn)=A_E(n)A_E(m)
 \qquad(m,n\geq1).
 \label{eq:compression-multiplicative}
\end{equation}
\end{proposition}

\begin{proof}
For a ramified label \(n\) we have \(\Psi_n\in A\) and \(E\in Z(A)\), so
\(E\Psi_n=\Psi_nE\).  For a good label \(n\) we have \(\Psi_n\in A'\) by
\cref{prop:good-in-commutant} and \(E\in A\), so again
\(E\Psi_n=\Psi_nE\).  A general \(n\) is a product of the two kinds by
\cref{prop:label-splitting}.  Therefore
\[
 E\Psi_mE\cdot E\Psi_nE=E^3\Psi_m\Psi_n=E\Psi_{mn},
\]
and restricting to \(V_E\) gives \eqref{eq:compression-multiplicative}; the
symmetry in \(m,n\) holds because \(\Psi_m\Psi_n=\Psi_{mn}=\Psi_n\Psi_m\).
\end{proof}

\section{The sporadic simple groups}
\label{sec:sporadic}

\subsection{Rational factors for the sporadic simple groups}

Names and class counts are those of the \emph{Atlas}~\cite{Atlas}.
Write
\[
 K_{-3}=\Q(\sqrt{-3}),\qquad K_5=\Q(\zeta_5).
\]
The degree of $K_5/\Q$ is four.  Hence a factor $M_r(K_5)$ gives four
conjugate copies of $M_r(\C)$ after scalar extension.

\begin{lemma}
\label{lem:schur-index}
Let $B=M_r(D)$ be a simple rational factor with center $K$, Schur index
$s=\sqrt{\dim_KD}$, and split degree $d=rs$.  Suppose the associated
$B$-module in rational class-coordinate space has complex multiplicity $m$.
Then $s$ divides both $d$ and $m$.  In particular, $\gcd(d,m)=1$ forces
$D=K$.  For the possible indices in the rational group-algebra setting
see~\cite{Yamada}.
\end{lemma}

\begin{proof}
For each embedding $K\hookrightarrow\C$, scalar extension sends $B$ to
$M_{rs}(\C)$.  A simple left $B$-module has $K$-dimension $rs^2$ and
complexifies to $s$ copies of the natural $rs$-dimensional module
\cite[\S{}29]{Reiner}.  Every
rational $B$-module therefore gives a complex multiplicity divisible by
$s$, while $s$ also divides $d=rs$.
\end{proof}

\begin{theorem}
\label{thm:sporadic-census}
The rational algebras of the twenty-six sporadic simple
groups are listed below.  The columns $z$ and $c$ give $\dim_{\Q}Z(A)$ and
$\dim_{\Q}A'$.  Every simple factor has Schur index one.  The only center
fields are $\Q$, $K_{-3}$, and $K_5$, and $K_5$ occurs only for $Ly$.

\begingroup
\footnotesize
\normalfont
\setlength{\tabcolsep}{3pt}
\renewcommand{\arraystretch}{1.08}
\begin{longtable}{@{}lrrrrp{0.54\textwidth}@{}}
\toprule
$G$ & $k(G)$ & $\dim A$ & $z$ & $c$ & $A$ over $\Q$\\
\midrule
\endfirsthead
\toprule
$G$ & $k(G)$ & $\dim A$ & $z$ & $c$ & $A$ over $\Q$\\
\midrule
\endhead
$M_{11}$ & 10 & 66 & 3 & 3 & $M_8(\Q)\oplus\Q^2$\\
$M_{12}$ & 15 & 149 & 3 & 3 & $M_{12}(\Q)\oplus M_2(\Q)\oplus\Q$\\
$M_{22}$ & 12 & 102 & 3 & 3 & $M_{10}(\Q)\oplus\Q^2$\\
$M_{23}$ & 17 & 151 & 5 & 5 & $M_{12}(\Q)\oplus M_2(\Q)\oplus\Q^3$\\
$M_{24}$ & 26 & 452 & 4 & 4 & $M_{21}(\Q)\oplus M_3(\Q)\oplus\Q^2$\\
$J_1$ & 15 & 111 & 4 & 4 & $M_{10}(\Q)\oplus M_3(\Q)\oplus K_{-3}$\\
$J_2$ & 21 & 269 & 3 & 3 & $M_{16}(\Q)\oplus M_3(\Q)\oplus M_2(\Q)$\\
$J_3$ & 21 & 209 & 6 & 6 & $M_{14}(\Q)\oplus M_3(\Q)\oplus\Q^2\oplus K_{-3}$\\
$J_4$ & 62 & 2000 & 11 & 11 &
$M_{44}(\Q)\oplus M_7(\Q)\oplus M_2(\Q)^2\oplus\Q\oplus K_{-3}^3$\\
$HS$ & 24 & 444 & 4 & 4 & $M_{21}(\Q)\oplus\Q^3$\\
$McL$ & 24 & 334 & 5 & 5 & $M_{18}(\Q)\oplus M_2(\Q)^2\oplus\Q^2$\\
$He$ & 33 & 647 & 5 & 5 & $M_{25}(\Q)\oplus M_4(\Q)\oplus M_2(\Q)\oplus\Q^2$\\
$Ru$ & 36 & 792 & 9 & 9 & $M_{28}(\Q)\oplus\Q^4\oplus K_{-3}^2$\\
$Suz$ & 43 & 1381 & 5 & 5 & $M_{37}(\Q)\oplus M_3(\Q)\oplus\Q^3$\\
$O'N$ & 30 & 411 & 9 & 12 & $M_{20}(\Q)\oplus M_2(\Q)\oplus\Q^5\oplus K_{-3}$\\
$Co_1$ & 101 & 9415 & 4 & 4 & $M_{97}(\Q)\oplus M_2(\Q)\oplus\Q^2$\\
$Co_2$ & 60 & 2926 & 5 & 5 & $M_{54}(\Q)\oplus M_2(\Q)^2\oplus\Q^2$\\
$Co_3$ & 42 & 1306 & 5 & 5 & $M_{36}(\Q)\oplus M_2(\Q)^2\oplus\Q^2$\\
$Fi_{22}$ & 65 & 3261 & 7 & 7 & $M_{57}(\Q)\oplus M_2(\Q)^2\oplus\Q^4$\\
$Fi_{23}$ & 98 & 8294 & 6 & 6 & $M_{91}(\Q)\oplus M_3(\Q)\oplus\Q^4$\\
$Fi_{24}'$ & 108 & 8853 & 12 & 15 & $M_{94}(\Q)\oplus M_2(\Q)^2\oplus\Q^9$\\
$HN$ & 54 & 1879 & 7 & 10 & $M_{43}(\Q)\oplus M_5(\Q)\oplus\Q^5$\\
$Ly$ & 53 & 1543 & 12 & 12 &
$M_{39}(\Q)\oplus M_3(\Q)\oplus M_2(\Q)\oplus\Q^3\oplus K_{-3}\oplus K_5$\\
$Th$ & 48 & 1384 & 10 & 10 & $M_{37}(\Q)\oplus M_2(\Q)^2\oplus\Q^7$\\
$\B$ & 184 & 26598 & 13 & 19 &
$M_{163}(\Q)\oplus M_3(\Q)\oplus M_2(\Q)^3\oplus\Q^8$\\
$\M$ & 194 & 28958 & 15 & 15 &
$M_{170}(\Q)\oplus M_5(\Q)\oplus M_3(\Q)\oplus M_2(\Q)^4\oplus\Q^8$\\
\bottomrule
\end{longtable}
\endgroup
\end{theorem}

\begin{proof}
The computation described in \cref{sec:certification} constructs each
rational word algebra and verifies closure.  A separating central element
gives square-free center factors; rational central idempotents give the
dimensions of the corresponding rational ideals.  These data determine the
split degree and class-space multiplicity of each complex block.  Their
gcd is one, so \cref{lem:schur-index} gives Schur
index one.

The quadratic defining polynomials have square-free discriminant $-3$.  The
only quartic center is defined by
$x^4+x^3+x^2+x+1$, which gives $K_5$.  Ranks at the primes
$1{,}000{,}003$ and $1{,}000{,}033$ reproduce the algebra, center, and
commutant dimensions of each residual calculation.
\end{proof}

\begin{corollary}
\label{cor:sporadic-saturation}
For each of the twenty-two sporadic groups whose class-coordinate
representation is multiplicity free, \(\Pram(G)=\PAd(G)\).  In particular the
Monster satisfies \(\Pram(\M)=\PAd(\M)\).
\end{corollary}

\begin{proof}
Immediate from \cref{thm:good-label-saturation} and the equality \(z=c\)
recorded in the table.
\end{proof}


The class-coordinate representation is multiplicity-free for 22 sporadic
groups.  The exceptions are $O'N$, $Fi_{24}'$, $HN$, and $\B$.
For the first three, a scalar complex block occurs with multiplicity two.
For $\B$, one scalar block and one $3$-dimensional block each occur
with multiplicity two.  The distinction between $z$ and $c$ in the table
makes these multiplicities visible.

\subsection{Good-label characters}
\label{sec:characters}

The sporadic census contains nine blocks with center \(\Q(\zeta_3)\) and
one block with center \(\Q(\zeta_5)\).

Throughout, \(N=\exp(G)\), \(U_N=(\Z/N)^\times\), \(\Gamma=\Psi(U_N)\), and
\(\chi_D\) denotes the Kronecker symbol \(\bigl(\tfrac D\cdot\bigr)\) of a
fundamental discriminant \(D\).

\subsubsection{Recovering the unit action from ramified data}

The CTblLib input consists of class sizes, element orders, and the power
maps \(\sigma_p\) for primes \(p\mid N\).

\begin{definition}
\label{def:witness-subgroup}
For an integer \(m\mid N\) put
\begin{equation}
 H_m=\bigl\langle\,p\bmod m\ :\ p\ \text{prime},\ p\mid N,\ p\nmid m\,\bigr\rangle
 \ \leq\ U_m .
 \label{eq:witness-subgroup}
\end{equation}
We call \(G\) \emph{ramified transparent} if \(H_m=U_m\) for every element
order \(m\) of \(G\).
\end{definition}

\begin{theorem}
\label{thm:unit-reconstruction}
Let \(C\) be a conjugacy class whose elements have order \(m\), let
\(u\in U_N\), and let \(v\) be any ramified label with \(v\equiv u\pmod m\).
Then
\begin{equation}
 \sigma_u(C)=\sigma_v(C).
 \label{eq:unit-reconstruction}
\end{equation}
If \(G\) is ramified transparent, the maps
\(\{\sigma_p:p\mid N\}\) determine the action of \(U_N\) on \(\Conj(G)\).
\end{theorem}

\begin{proof}
For \(x\in C\) we have \(x^m=1\), so \(x^u=x^{u\bmod m}=x^{v\bmod m}=x^v\),
whence \(\sigma_u(C)=\sigma_v(C)\), which is
\eqref{eq:unit-reconstruction}.  A ramified \(v\) with \(v\equiv u\pmod m\)
exists if and only if \(u\bmod m\in H_m\), because the residues of ramified
labels modulo \(m\) that are units are the elements of \(H_m\):
a ramified \(v\) is a product of primes dividing \(N\), those dividing \(m\)
contribute non-units, and \(\gcd(v,m)=1\) forces every prime factor of \(v\)
to be coprime to \(m\).  Under ramified transparency this holds for all
\(m\) and \(u\), and \(\sigma_v\) is computed from the stored maps by
multiplicativity, \(\sigma_v=\prod_p\sigma_p^{e_p}\) for \(v=\prod_pp^{e_p}\).
\end{proof}

\begin{proposition}
\label{prop:reconstruction-coverage}
Twenty-four of the twenty-six sporadic simple groups are ramified
transparent.  The exceptions are \(J_3\), where
\([U_{15}:H_{15}]=2\), and \(Th\), where \([U_{30}:H_{30}]=2\); in both cases
the missing coset is the one containing \(-1\), and the ambiguity concerns a
single pair of classes.
\end{proposition}

\begin{proof}
Closing the residues of the primes dividing \(N\) under multiplication
modulo \(m\) and comparing \(|H_m|\) with \(\varphi(m)\) leaves only the two
displayed failures.  For \(J_3\), \(N=2^3\cdot3^2\cdot
5\cdot17\cdot19\) and the primes coprime to \(15\) reduce to
\(2,17,19\equiv2,2,4\pmod{15}\), generating
\(\langle2\rangle=\{1,2,4,8\}\) of index two in \(U_{15}\); the missing coset
is \(7\langle2\rangle=\{7,11,13,14\}\ni-1\).  The computation for \(Th\) at
\(m=30\) is identical in form.
\end{proof}

For \(J_3\) and \(Th\), call a permutation of \(\Conj(G)\)
\emph{admissible} if it extends the forced values, preserves class size and
element order, and commutes with the stored power maps.

\begin{proposition}
\label{prop:ambiguity-harmless}
For each sporadic simple group and each chosen generator of \(U_N\), the
admissible completion is unique.  The stored ramified maps, class sizes, and
element orders determine the action of \(U_N\) on \(\Conj(G)\).
\end{proposition}

\begin{proof}
For the twenty-four ramified transparent groups,
\cref{thm:unit-reconstruction} determines all images.  For \(J_3\), two of
the six chosen generators, those with nontrivial \(3\)- and
\(5\)-components, leave undetermined the two
classes of element order \(15\), the remaining nineteen images being forced;
of the two possible assignments only one commutes with \(\sigma_3\).  For
\(Th\), one of the eight generators, the \(3\)-component, leaves the
two classes of element order \(30\) undetermined out of forty-eight, and
there the disambiguating constraint is \(\sigma_2\).  Which of the two
assignments survives is not uniform: the surviving completion fixes both
classes for the \(3\)-component of \(J_3\) and transposes them in the other
two cases.
\end{proof}

The statement is independent of the displayed generators: once the
permutations are recovered on any generating set, the homomorphism
\(U_N\to\operatorname{Sym}(\Conj(G))\) determines every other choice.

\subsubsection{The character census}

\begin{theorem}
\label{thm:good-label-census}
Let \(G\) be a sporadic simple group and \(E_k\) a primitive central
idempotent of \(\Pram(G)\).  There is a character
\(\lambda_k:U_N\to\mu(K_k)\) with
\[
 \Psi_uE_k=\lambda_k(u)E_k .
\]
The principal block carries the trivial character.  Of the remaining blocks,
ten have nonrational center and are listed in
\cref{sec:nonrational-blocks}; every other one carries the trivial character
or the Kronecker symbol \(\chi_{D_k}\) of a fundamental discriminant \(D_k\)
whose prime divisors divide \(N\).  Five blocks, in \(O'N\), \(Fi_{24}'\),
\(HN\), and \(\B\), have multiplicity two, and \(\Psi_u\) is scalar on those
as well.
\end{theorem}

\begin{proof}
Compute the primitive central idempotents over \(\Q\) and the operators
\(\Psi_u\) for one generator per cyclic factor of
\(\prod_{p^a\|N}U_{p^a}\).  The block-local theorem treats the
multiplicity-one blocks.  Rational row reduction gives
\(\Psi_uE_k=\lambda_k(u)E_k\) on the five multiplicity-two blocks.  The
order of \(\lambda_k(u)\) is the least \(e\) with
\((\Psi_uE_k)^e=E_k\); the character
order is the least common multiple over the generators; the conductor is
assembled from the local orders by the standard rule for
\(U_{p^a}\).  For a rational block the parity \(\lambda_k(-1)\) selects the
sign of the fundamental discriminant of the given conductor.  Direct
substitution gives \(\lambda_k(u)=\chi_D(u)\) on the generators.

The principal block character is trivial.  That block is the cyclic
module \(A\delta_{1A}\) generated by the identity-class coordinate
(\cref{sec:certification}); every \(\Psi_u\) fixes \(\delta_{1A}\) and
commutes with \(A\), so it fixes \(A\delta_{1A}\) pointwise.
\end{proof}

\subsubsection{The ten nonrational blocks}
\label{sec:nonrational-blocks}

The ten nonrational blocks occur in
\[
 J_1,\qquad J_3,\qquad J_4,\qquad Ru,\qquad O'N,\qquad Ly.
\]
Nine blocks have center \(\Q(\zeta_3)\); the unique remaining block has
center \(\Q(\zeta_5)\) and belongs to \(Ly\).  In every case the star acts as
complex conjugation on the center, the character order is the odd prime
\(e\) with \(K=\Q(\zeta_e)\), and the block is supported on \(e\)
classes of a single element order.  Its split degree is
\(d_k=\varphi(e)=e-1\).  The complete list is:

\begin{center}
\normalfont
\begin{tabular}{@{}llrrr@{}}
\toprule
$G$ & $K_k$ & $e$ & cond. & elt.\ order\\
\midrule
$J_1$  & $K_{-3}$ & 3 & 19 & 19\\
$J_3$  & $K_{-3}$ & 3 &  9 &  9\\
$J_4$  & $K_{-3}$ & 3 & 37 & 37\\
$J_4$  & $K_{-3}$ & 3 & 31 & 31\\
$J_4$  & $K_{-3}$ & 3 & 43 & 43\\
$Ru$   & $K_{-3}$ & 3 &  7 & 14\\
$Ru$   & $K_{-3}$ & 3 & 13 & 26\\
$O'N$  & $K_{-3}$ & 3 & 19 & 19\\
$Ly$   & $K_{-3}$ & 3 & 67 & 67\\
$Ly$   & $K_5$    & 5 & 31 & 31\\
\bottomrule
\end{tabular}
\end{center}


\begin{remark}
\label{rem:pariahs}
The six groups carrying a nonrational center are exactly the six pariah
groups, the sporadic groups not involved in the Monster.  Every group of
the happy family has \(\Pram(G)\) split over \(\Q\) with all centers
rational.  We record this as an observation of the census; no explanation
is offered here.
\end{remark}


\begin{proposition}
\label{prop:conductor-determined}
Let \(E\) be a block whose support consists of classes of a common element
order \(m\), and let \(\lambda_E\) be its good-label character, of order
\(e>1\).  Write \(m=2^am'\) with \(m'\) odd.  Then
\begin{enumerate}
\item \(\lambda_E\) factors through \(U_m\), so its conductor divides \(m\);
\item if \(a\leq1\), the conductor divides \(m'\);
\item if moreover \(m'=p^b\) is a prime power and
\(e\nmid\varphi(p^{b-1})\), the conductor equals \(m'\).
\end{enumerate}
\end{proposition}

\begin{proof}
(1) By \cref{thm:unit-reconstruction} the permutation \(\sigma_u\) acts on a
class of element order \(m\) through \(u\bmod m\) alone, so \(\lambda_E\)
factors through \(U_m\) and its conductor divides \(m\).

(2) For odd \(m'\) the Chinese remainder theorem gives
\(U_{2m'}\cong U_2\times U_{m'}=U_{m'}\), so when \(a\leq1\) the character
factors through \(U_{m'}\).

(3) The conductor is \(p^c\) for some \(c\leq b\), minimal such that
\(\lambda_E\) factors through \(U_{p^c}\); then
\(e=|\lambda_E(U_N)|\) divides \(|U_{p^c}|=\varphi(p^c)\).  For \(p\) odd and
\(c\leq b-1\) one has \(\varphi(p^c)\mid\varphi(p^{b-1})\), whence
\(e\mid\varphi(p^{b-1})\), contrary to hypothesis.  Hence \(c=b\).
\end{proof}

\begin{corollary}
\label{cor:conductor-odd-part}
In every row of \cref{sec:nonrational-blocks} the conductor equals the odd
part \(m'\) of the common element order.
\end{corollary}

\begin{proof}
Each row has \(a\leq1\), and \(m'\) is a prime power.  In nine rows \(m'\) is
one of \(7,13,19,31,37,43,67\); then \(b=1\), and \(e>1\) does not divide
\(\varphi(p^0)=1\).  In the remaining row \(m'=9\) and \(e=3\), which does
not divide \(\varphi(3)=2\).
\Cref{prop:conductor-determined} applies in every case.
\end{proof}

\begin{example}
\label{ex:ly-quintic}
Among the twenty-six sporadic groups, \(Ly\) alone has a quartic center.
Its five conjugacy classes of elements of order \(31\) form one rational
class and carry the full \(\Q(\zeta_5)\) factor of \(\Pram(Ly)\).  Writing
\(R\) for this five-class orbit, the block is
\[
 V_k=\Bigl\{f\in\Span_\Q(\delta_C:C\in R)\ :\ \textstyle\sum_{C\in R}f(C)=0\Bigr\},
 \qquad \dim_\Q V_k=4,
\]
and the good labels act through a quintic character of conductor \(31\).
Thus \(V_k\) is one-dimensional over its cyclotomic endomorphism field:
\[
 \End_{\Q[U_N]}(V_k)=\Q(\zeta_5).
\]
The character is even, so inversion fixes all five order-\(31\) classes.
\end{example}

\begin{example}
\label{ex:j1-cubic}
The three classes of elements of order \(19\) in \(J_1\) form a single
rational class; the augmentation-zero part of their span is a two-dimensional
block with center \(\Q(\sqrt{-3})=\Q(\zeta_3)\), on which the good labels act
through a cubic character of conductor \(19\).  Since \(19\equiv1\pmod3\),
such characters exist modulo \(19\); there are two of them, complex conjugate
to one another, and they are interchanged by the choice of embedding
\(\Q(\zeta_3)\hookrightarrow\C\).  The pair, not either member, is the
invariant.
\end{example}

\begin{theorem}
\label{thm:regular-rational-class}
Let \(R\subseteq\Conj(G)\) be a \(\Gamma\)-orbit with \(|R|=r\), let
\[
 W_R=\Bigl\{f\in\Span_\Q(\delta_C:C\in R):\textstyle\sum_{C\in R}f(C)=0\Bigr\},
 \qquad \dim_\Q W_R=r-1 .
\]
Then \(\Gamma\) acts regularly on \(R\) through an abelian quotient
\(A_R\) of order \(r\), and \(W_R\) is \(\Gamma\)-stable.  If \(r\) is
prime, then \(A_R\cong\Z/r\), and \(\Gamma\) acts on \(W_R\) through a
character of order \(r\) with values in \(\Q(\zeta_r)\), and
\(\End_{\Q[\Gamma]}(W_R)=\Q(\zeta_r)\).  If moreover \(W_R\) is a single
split block of \(\Pram(G)\) of multiplicity one, that block has center
\(K=\Q(\zeta_r)\) and good-label character of order \(r\).
\end{theorem}

\begin{proof}
\(\Gamma\) is abelian by \cref{lem:good-labels}(4) and acts transitively on
\(R\) by definition of an orbit.  In a transitive action of an abelian group
all point stabilizers agree; after quotienting by their common value
\(\Gamma_R\), the action is free and transitive.  Thus
\(A_R=\Gamma/\Gamma_R\) has order \(r\), and \(R\) is an \(A_R\)-torsor.
If \(r\) is prime then \(A_R\cong\Z/r\).  The span of \(R\) is the regular
representation \(\Q[A_R]\), which splits as \(\Q\oplus W_R\) with \(\Q\) the
invariants; both summands are \(\Gamma\)-stable.

Suppose \(r\) is prime.  Then \(\Q[\Z/r]\cong\Q\times\Q(\zeta_r)\)~\cite{Washington}, the
second factor acting on \(W_R\), so \(W_R\) is a one-dimensional
\(\Q(\zeta_r)\)-vector space on which a generator of \(\Z/r\) acts as
multiplication by \(\zeta_r\).  Hence the induced character has order
\(r\), its values generate \(\Q(\zeta_r)\), and the commuting algebra is
\(\Q(\zeta_r)\).  If \(W_R\) is a single split multiplicity-one block of
\(\Pram(G)\), its center is its commuting algebra inside the isotypic
component, which contains and is contained in \(\Q(\zeta_r)\) by the previous
sentence and \cref{thm:block-local-character}.
\end{proof}

\begin{corollary}
\label{cor:orbit-divisibility}
Let \(r\) be an odd prime.  Then \(V_G\) has a \(\Q[U_N]\)-summand with
endomorphism field \(\Q(\zeta_r)\) if and only if \(r\) divides \(|R|\) for
some rational class \(R\).
\end{corollary}

\begin{proof}
By \cref{thm:rational-class-decomposition},
\(V_G\cong\bigoplus_R\Q[A_R]\) with \(A_R=U_N/\Gamma_R\) abelian of order
\(|R|\).  For an abelian group \(A\), the simple summands of \(\Q[A]\) are the
fields \(\Q(\chi)=\Q(\zeta_{\mathrm{ord}\chi})\) as \(\chi\) runs over the
Galois orbits of characters of \(A\), each acting on itself, so the
endomorphism field of that summand is \(\Q(\chi)\).  Since \(r\) is an odd
prime, \(\Q(\chi)=\Q(\zeta_r)\) holds if and only if
\(\mathrm{ord}\,\chi\) is
\(r\) or \(2r\), and such a \(\chi\) exists for some \(A_R\) if and only if
\(r\) divides the exponent of some \(A_R\), equivalently \(r\mid|A_R|=|R|\).
\end{proof}

\begin{corollary}
\label{cor:nonrational-from-orbits}
Every nonrational block in the sporadic census arises as in
\cref{thm:regular-rational-class}: nine from rational classes of size three,
giving \(K=\Q(\zeta_3)=\Q(\sqrt{-3})\) and cubic characters, and one, the
\(\Q(\zeta_5)\) block of \(Ly\), from a rational class of size five, giving a
quintic character.  In every case \(d_k=r-1\).
\end{corollary}

\begin{proof}
For each block, compute its class support.  For all ten nonrational blocks the
support is a single orbit \(R\) of prime size \(r\in\{3,5\}\), all of whose
classes have one element order, and \(d_k=|R|-1\); the block therefore is
\(W_R\) and \cref{thm:regular-rational-class} applies.  Comparing its
conclusion with the computed columns, \(K_k=\Q(\zeta_r)\) and the character
order is \(r\), as tabulated.
\end{proof}

\Cref{thm:regular-rational-class} with \(r=2\) accounts for a large part of
the rational half of \cref{thm:good-label-census} as well: a rational class
splitting into two conjugacy classes contributes a one-dimensional block on
which \(\Gamma\) acts by a quadratic character.  The blocks with
\(d_k>r-1\) are the ones where several rational classes fuse into a single
simple factor; those blocks have a
conductor divisible by more than one prime.

\subsubsection{The Monster}

\begin{proposition}
\label{prop:monster-imaginary}
Let \(\M\) be the Monster, so that \(k(\M)=194\) and
\begin{equation}
 N=\exp(\M)=2^5\,3^3\,5^2\cdot7\cdot11\cdot13\cdot17\cdot19\cdot23
 \cdot29\cdot31\cdot41\cdot47\cdot59\cdot71 .
 \label{eq:monster-exponent}
\end{equation}
Its ramified algebra is
\[
 \Pram(\M)\cong
 M_{170}(\Q)\oplus M_5(\Q)\oplus M_3(\Q)
 \oplus M_2(\Q)^4\oplus\Q^8.
\]
Thus all fifteen blocks have center \(\Q\).  Two carry the trivial character;
the other thirteen carry quadratic characters whose
fundamental discriminants are
\begin{equation}
 -11,\,-20,\,-23,\,-31,\,-39,\,-47,\,-56,\,-59,\,-71,\,-87,\,-95,\,-104,\,-119 .
 \label{eq:monster-discriminants}
\end{equation}
All thirteen are negative, so every nontrivial good-label character of the
Monster is odd: \(\lambda_k(-1)=-1\) on each of the thirteen blocks.
\end{proposition}

\begin{proof}
The rational decomposition and the character conductors are the Monster row
of \cref{thm:sporadic-census,thm:good-label-census}.  The inversion map in
the CTblLib class ordering fixes \(172\) of the \(194\) classes, while the
unit-group action has \(172\) orbits.  Hence the numbers of real and rational
classes agree.  The defining module is multiplicity free, so
\cref{prop:real-rational-parity} forces every nontrivial block character to
be odd.  The thirteen fundamental discriminants are then determined by
their certified conductors and their values on the unit-group generators.
\end{proof}

By \cref{prop:rational-splitting}, the
\(+1\) eigenspace of \(\Psi_{-1}\) is spanned by the orbit sums of the
inversion action, so \eqref{eq:monster-discriminants} says that
every nontrivially twisted block of the Monster is detected already by
inversion: the twenty-two dimensions \(194-172\) lie entirely in the
\(-1\) eigenspace of complex conjugation on classes.  This is special to \(\mathbb M\) among the large
sporadic groups: \(Ly\), for instance, has four blocks with positive
discriminant \(40,37,24,21\) and one with \(-11\), and \(J_4\) has three
positive and one negative.

The assignment
\[
 G\ \longmapsto\ \bigl\{(K_k,\ d_k,\ \lambda_k)\bigr\}_{k}
\]
is an invariant of the marked weighted power system.  The groups \(Co_2\)
and \(Co_3\) have
the same rational block list \(M_r(\Q)\oplus M_2(\Q)^2\oplus\Q^2\) up to the
principal degree, but their residual characters differ, \(\chi_{-23},\chi_{-7}\)
and \(\chi_{-20},\chi_{-23}\), respectively.


\subsection{Rank-one saturation}
\label{sec:rank-one-saturation}

Let a finite elementary abelian \(2\)-group \(K\) act by weight-preserving
permutations on a finite weighted set \(X\).  Write
\[
 V_X=\bigoplus_{\chi\in\widehat K}V_\chi,
 \qquad
 V_\chi=\{v:kv=\chi(k)v\text{ for all }k\in K\}.
\]
Every \(V_\chi\) is rational.  Hence
\[
 \End_{\Q[K]}(V_X)=\bigoplus_{\chi:m_\chi>0}M_{m_\chi}(\Q),
 \qquad m_\chi=\dim_\Q V_\chi.
\]

\begin{lemma}[Cyclic rank-one generation]
\label{lem:cyclic-rank-one}
Let \(W\) be a finite-dimensional rational inner-product space, let
\(B\subseteq\End_\Q(W)\) be a unital star-algebra, and let \(v\in W\).
If the orthogonal rank-one projection \(E_v\) belongs to \(B\) and
\(Bv=W\), then \(B=\End_\Q(W)\).
\end{lemma}

\begin{proof}
Choose \(b_1,\ldots,b_d\in B\) such that \(b_1v,\ldots,b_dv\) is a basis
of \(W\).  The operators \(b_iE_vb_j^*\) are nonzero scalar multiples of
the matrix units for this basis.  They span \(\End_\Q(W)\).
\end{proof}

For \(a\) in the monoid generated by the labels, put
\[
 r_a=P_a^*P_a\mathbf1.
\]
The transfer identity makes \(P_a^*P_a=M_{r_a}\).  A finite tuple of the
functions \(r_a\) will be called a root profile.

\begin{lemma}[Root-profile projectors]
\label{lem:root-profile-projectors}
Let \(O\subseteq X\) be a \(K\)-orbit.  If a finite root profile takes one
joint value on \(O\) and a different joint value on every other orbit,
then the coordinate projection \(E_O\) belongs to \(\PE(\mathscr X)\).
\end{lemma}

\begin{proof}
The diagonal operators \(M_{r_a}\) belong to \(\PE(\mathscr X)\).
Multivariate Lagrange interpolation in the finitely many joint values
produces the indicator of \(O\).
\end{proof}

Suppose \(O=\{x,y\}\) is a two-point orbit belonging to the character
\(\chi\).  Since \(K\) preserves the weight, \(\mu(x)=\mu(y)\), and
\[
 E_O=E_{\delta_x+\delta_y}+E_{\delta_x-\delta_y}
\]
after normalizing the two orthogonal vectors.  Once the trivial matrix
block has been generated, an isolated root profile for \(O\) therefore
supplies the rank-one projection onto
\(\Q(\delta_x-\delta_y)\subseteq V_\chi\).

\begin{theorem}[Root-cyclic saturation]
\label{thm:root-cyclic-saturation}
Let \(A=\PE(\mathscr X)\subseteq\End_{\Q[K]}(V_X)\), where \(K\) is
elementary abelian of exponent \(2\).  Assume:
\begin{enumerate}
\item a fixed point \(x_0\) has \(E_{x_0}\in A\), and
      \(A\delta_{x_0}=V_{\mathbf1}\);
\item for each nontrivial \(\chi\) with \(m_\chi>1\), a two-point orbit
      \(O_\chi=\{x_\chi,y_\chi\}\) has an isolated root profile and
      \(A(\delta_{x_\chi}-\delta_{y_\chi})=V_\chi\);
\item the restrictions of elements of \(A\) separate the characters
      \(\chi\) with \(m_\chi=1\).
\end{enumerate}
Then
\[
 A=\End_{\Q[K]}(V_X).
\]
The saturation certificate requires
\[
 m_{\mathbf1}+
 \sum_{\substack{\chi\ne\mathbf1\\m_\chi>1}}m_\chi
\]
orbit vectors, together with scalar eigenvalue data on the
one-dimensional channels.
\end{theorem}

\begin{proof}
The first hypothesis and \cref{lem:cyclic-rank-one} generate the full
trivial block.  For every higher-dimensional nontrivial channel,
\cref{lem:root-profile-projectors} and subtraction of the trivial
rank-one summand give the projection onto
\(\Q(\delta_{x_\chi}-\delta_{y_\chi})\).  The cyclic rank-one lemma then
generates \(\End_\Q(V_\chi)\).  Project away these matrix blocks.  On the
remaining one-dimensional channels, the separating restrictions generate
all coordinate idempotents by interpolation.  Thus \(A\) contains each
simple factor of the displayed commutant.
\end{proof}

\begin{lemma}[Modular cyclicity]
\label{lem:modular-cyclicity}
Let \(R=\Z[1/M]\), let \(B\subseteq M_d(R)\) be generated by fixed
matrices, and let \(v\in R^d\).  Let \(\ell\nmid M\) be prime.  If words
\(w_1,\ldots,w_d\) in the generators satisfy
\[
 \det[w_1v\ \cdots\ w_dv]\not\equiv0\pmod\ell,
\]
then \(v\) is cyclic over \(\Q\).
\end{lemma}

\begin{proof}
The orbit determinant belongs to \(R\).  Its nonzero reduction cannot
come from the zero rational number.
\end{proof}

\subsection{The Monster symmetry commutant}

Let \(\Gamma_{\M}\) be the group of permutations of the \(194\) Monster
classes that preserve class size and commute with every prime-power map.  It
is the automorphism group of the marked weighted power system.

\begin{theorem}[Weighted-power symmetry and commutant]
\label{thm:monster-symmetry-commutant}
The group \(\Gamma_{\M}\) is elementary abelian of rank \(14\).  Its
subgroup generated by the nineteen stored ordinary-table automorphisms has
rank \(13\), and a complementary factor is generated by
the involution
\[
 (16B\ 16C)(32A\ 32B).
\]
For the permutation representation on
\(V_{\M}=\Q^{\operatorname{Conj}(\M)}\),
\[
 \operatorname{End}_{\Q\Gamma_{\M}}(V_{\M})
 =\Pram(\M)
 \cong
 M_{170}(\Q)\oplus M_5(\Q)\oplus M_3(\Q)
 \oplus M_2(\Q)^{\oplus4}\oplus\Q^{\oplus8}.
\]
Its dimension is \(28958\); its center and commutant have dimension \(15\).
\end{theorem}

\begin{proof}
Start with the coloring by class size and element order.  Refine it by the
colors of all prime-power images and by the weighted color counts in each
reverse fibre.  The first two refinement rounds have \(161\) and \(170\)
cells, and the second coloring is stable.  Its nonsingleton cells are
twenty-four pairs.  Commutation with the fifteen prime-power maps joins
their swap variables into fourteen components of sizes
\[
 2,5,3,2,1,2,2,1,1,1,1,1,1,1.
\]
Every component swap satisfies all power-map equations.  Hence
\(\Gamma_{\M}\cong(C_2)^{14}\).  The binary matrix of the nineteen stored
table automorphisms has rank \(13\); its missing coordinate is the component
\((16B\ 16C)(32A\ 32B)\).

Each singleton contributes one trivial vector.  Each pair contributes its
sum to the trivial channel and its difference to the character of its swap
component.  Thus
\[
 V_{\M}\cong
 \Q_{\mathbf1}^{170}\oplus
 \Q_{\chi_1}^{2}\oplus\Q_{\chi_2}^{5}\oplus
 \Q_{\chi_3}^{3}\oplus\Q_{\chi_4}^{2}\oplus
 \Q_{\chi_6}^{2}\oplus\Q_{\chi_7}^{2}
 \oplus\bigoplus_{i\in\{5,8,9,10,11,12,13,14\}}\Q_{\chi_i}.
\]
This gives the displayed target commutant.

It remains to prove saturation.  Since every Monster element has exponent
dividing \(N=\exp(\M)\), the operator
\(|\M|^{-1}P_N^*P_N\) is the coordinate projection onto \(1A\).
Modulo \(1000003\), the orbit of \(\delta_{1A}\) has rank \(170\), with a
retained basis of words of length at most six.  The pairs
\[
 16B/C,\ 23A/B,\ 31A/B,\ 39C/D,\ 44A/B,\ 47A/B
\]
have isolated root profiles for the six nontrivial channels of dimension
greater than one.  Their orbit determinants modulo \(1000003\) are
\[
 1,\ 2,\ 1,\ 1,\ -1,\ 1,
\]
with maximum word lengths \(1,3,1,1,1,1\).  Finally, the restrictions of
\((P_2,P_5,P_7)\) to the eight scalar channels are pairwise distinct.
The hypotheses of \cref{thm:root-cyclic-saturation} hold, so
\(\Pram(\M)=\End_{\Q\Gamma_{\M}}(V_{\M})\).  The certificate uses
\(170+2+5+3+2+2+2=186\) retained orbit vectors.  All denominators are
invertible modulo \(1000003\), so \cref{lem:modular-cyclicity} lifts the
seven orbit determinants to \(\Q\).
\end{proof}

\begin{remark}
The exceptional rational \(M_2(\Q)\)-channel is the unique symmetry
direction on which the complementary involution acts and the stored
table-automorphism subgroup acts trivially.  No identification of the
stored subgroup with the full unit-power image is used.
\end{remark}



\subsection{Symmetries invisible to the character table}
\label{sec:invisible-symmetries}

Galois conjugation of classes is an automorphism of the character table
and of the weighted power system, so \(\Psi(U_N)\subseteq\Gamma_G\) for
every \(G\), where \(\Gamma_G\) is the automorphism group of the weighted
power system.  \Cref{thm:monster-symmetry-commutant} shows that the
inclusion can be strict: for the Monster, \(\Gamma_{\M}\cong(C_2)^{14}\)
while the Galois image has rank \(13\), the thirteen distinct quadratic
characters of \cref{prop:monster-imaginary}.  The complementary involution
\((16B\ 16C)(32A\ 32B)\) preserves class sizes and all fifteen prime-power
maps.  It lies outside the group generated by the stored automorphisms of the
character table, so it is not induced by any automorphism of that table.
Thus the classes \(16B,16C\) and \(32A,32B\) are distinguished by the
character table but not by class sizes and power maps.

At the level of the algebra this involution is visible as the unique
nonprincipal block of \(\Pram(\M)\) on which the good labels act
trivially: by \cref{prop:rational-splitting} the Galois-fixed space
\(V_{\M}^{\Gamma}\) has dimension \(172\), and it decomposes under
\(\Pram(\M)\) as the \(170\)-dimensional principal block plus one
\(M_2(\Q)\)-block, the difference channel of the complementary involution.

The same signature occurs elsewhere in the census.  In the character table
of \cref{thm:good-label-census}, \(Co_2\) and \(Co_3\) each have an
\(M_2(\Q)\)-block carrying the trivial character, so each has a
two-dimensional nonprincipal simple factor inside its Galois-fixed space.
By \cref{thm:balanced-commutant} that factor comes from a balanced kernel
commuting with \(\Psi(U_N)\) and not accounted for by it.  Whether it
arises from a permutation of classes, as for the Monster, is not
determined by the computations reported here; the symmetry group
\(\Gamma_G\) was computed only for \(\M\).

For \(Fi_{23}\) the answer is known and is negative.  Its census row has a
scalar block with trivial character spanned by
\[
 v=\delta_{12H}-\tfrac13\,\delta_{12L}-\tfrac14\,\delta_{12M}
   +\tfrac1{12}\,\delta_{12O},
\]
a joint eigenvector of all sixteen generators, with eigenvalue \(0\) for
\(\Psi_2,\Psi_2^*,\Psi_3,\Psi_3^*\) and \(1\) for the others, orthogonal to
the principal block \(A\delta_{1A}\), which has exact dimension \(91\)
(\cref{sec:certification}).  The four coefficients of \(v\) have distinct
absolute values, so no permutation of classes commuting with the power maps
moves \(v\) except trivially; every element of \(\Q\Gamma_{Fi_{23}}\) acts
on the Galois-fixed space \(V^\Gamma\) as a scalar, while the projection
onto \(\Q v\) lies in the commutant and acts on \(V^\Gamma\) nontrivially.
Hence \(\Bal(Fi_{23})\supsetneq\Q\Gamma_{Fi_{23}}\): the commutant of
\(\Pram(Fi_{23})\) contains a balanced kernel that is not a combination of
permutation symmetries.  The Monster's extra block is a permutation
symmetry invisible to the character table; the \(Fi_{23}\) block is a
fractional symmetry invisible even to the automorphism group of the weighted
power system.

\subsection{Computational record}
\label{sec:characters-record}

The file \path{good_label_characters.json} records, per group, the generators
of \(U_N\) and the reconstruction data of \cref{thm:unit-reconstruction}, and
per primitive central idempotent the center field and its defining
polynomial, \(d_k\), the class support with its element orders, whether that
support is a single \(\Gamma\)-orbit, the values \(\lambda_k(u)\) at the
generators with their minimal polynomials and orders, and the conductor,
parity, and fundamental discriminant of \(\lambda_k\).  It checks
\eqref{eq:lefschetz}, \eqref{eq:trivial-block-sum},
\(\sigma_u\sigma_v=\sigma_{uv}\), and the admissible completions of
\cref{prop:ambiguity-harmless}; each generator has one completion.  Every
block acts by center scalars and every quadratic block agrees with its
Kronecker symbol at every generator.

\section{Infinite families and a comparison panel}
\label{sec:families}

\subsection{Groups of prime exponent}

A single many-to-one power map and its weighted reverse map already create
a noncommutative factor.

\begin{theorem}
\label{thm:prime-exponent-structural}
Suppose \(\exp(G)=p\) is prime and write \(k=|\operatorname{Conj}(G)|\).
Then
\begin{equation}
 \Pram(G)\cong
 \begin{cases}
 M_2(\Q),&k=2,\\
 M_2(\Q)\oplus\Q,&k>2.
 \end{cases}
 \label{eq:prime-exponent-type-structural}
\end{equation}
In the defining representation, the scalar factor in the second line has
multiplicity \(k-2\).
\end{theorem}

\begin{proof}
Let \(e=e_1\in\Q^k\), let \(u=(1,\ldots,1)^{\mathsf T}\), and let
\(c=(c_1,\ldots,c_k)^{\mathsf T}\) be the class-size vector, so \(c_1=1\) and
\(\sum_ic_i=|G|\).  Since \(x^p=1\) for every \(x\in G\), we have
\(\sigma_p(i)=1\) for all \(i\), so \eqref{eq:power-matrix} gives
\(P_p=ue^{\mathsf T}\); and
\(P_p^*=D_G^{-1}P_p^{\mathsf T}D_G
=D_G^{-1}eu^{\mathsf T}D_G=(|G|/c_1)\,e\,(c^{\mathsf T}/|G|)\), so
\begin{equation}
 P_p=ue^{\mathsf T},
 \qquad
 P_p^*=ec^{\mathsf T}.
 \label{eq:prime-rank-one-matrices}
\end{equation}

\emph{The decomposition.}  Put \(W=\Span_\Q\{e,u\}\) and
\[
 U=\{z\in\Q^k:z_1=0,\ c^{\mathsf T}z=0\}.
\]
Since \(\langle z,e\rangle_G=z_1/|G|\) and
\(\langle z,u\rangle_G=c^{\mathsf T}z/|G|\), the space \(U\) is the
class-orthogonal complement of \(W\).  If \(k\geq2\) then \(u\notin\Q e\), so
\(\dim W=2\); the two functionals \(z\mapsto z_1\) and
\(z\mapsto c^{\mathsf T}z\) are linearly independent because
\(c\neq c_1e=e\) when \(k\geq2\), so \(\dim U=k-2\).  Also
\(W\cap U=0\): if \(\alpha e+\beta u\in U\) then its first coordinate gives
\(\alpha+\beta=0\), while
\(c^{\mathsf T}(\alpha e+\beta u)=\alpha+\beta|G|=0\); subtracting and using
\(|G|>1\) forces \(\alpha=\beta=0\).  Hence
\(\Q^k=W\oplus U\) orthogonally.

Both generators vanish on \(U\): for \(z\in U\),
\(P_pz=u(e^{\mathsf T}z)=z_1u=0\) and
\(P_p^*z=e(c^{\mathsf T}z)=0\).  Both preserve \(W\), since
\(P_pe=u\), \(P_pu=u\), \(P_p^*e=c_1e=e\) and \(P_p^*u=(c^{\mathsf T}u)e=|G|e\).
So \(W\) and \(U\) are reducing for \(P_p\) and \(P_p^*\).

\emph{The action on \(W\).}  In the ordered basis \((e,u)\) the four
operators \(P_p,P_p^*,P_pP_p^*,P_p^*P_p\) restrict to
\begin{equation}
 \begin{aligned}
 P_p|_W&=
 \begin{pmatrix}0&0\\1&1\end{pmatrix},
 & P_p^*|_W&=
 \begin{pmatrix}1&|G|\\0&0\end{pmatrix},\\
 (P_pP_p^*)|_W&=
 \begin{pmatrix}0&0\\1&|G|\end{pmatrix},
 & (P_p^*P_p)|_W&=
 \begin{pmatrix}|G|&|G|\\0&0\end{pmatrix}.
 \end{aligned}
 \label{eq:prime-two-by-two}
\end{equation}
Reading these as vectors in \(\Q^4\) with coordinates
\((a_{11},a_{12},a_{21},a_{22})\) gives the matrix
\[
 \begin{pmatrix}
 0&0&1&1\\
 1&|G|&0&0\\
 0&0&1&|G|\\
 |G|&|G|&0&0
 \end{pmatrix},
\]
which is block anti-diagonal and has determinant
\(\pm|G|\,(|G|-1)^2\neq0\) because \(|G|>1\).  Hence the
four restrictions are linearly independent and span \(M_2(\Q)\).

\emph{Conclusion.}  Every word in \(P_p,P_p^*\) of length at least one
vanishes on \(U\), by the previous paragraph, and the span of such words
restricted to \(W\) is all of \(M_2(\Q)\) by the determinant computation.
Therefore
\[
 \PE(G)
 =\{a\oplus\lambda I_U:a\in\End_\Q(W),\ \lambda\in\Q\},
\]
the summand \(\lambda I_U\) arising only from the adjoined identity.  If
\(U=0\), that is \(k=2\), this is \(M_2(\Q)\).  If \(U\neq0\), the element
\(I_W\oplus0\) lies in the algebra, so the algebra is
\(M_2(\Q)\oplus\Q\), and the scalar factor acts on \(U\), which has
dimension \(k-2\).
\end{proof}

The theorem gives a family independent of the
classification of the underlying groups.  It applies to every elementary
abelian \(p\)-group and, for odd \(p\), to every extraspecial group of
exponent \(p\).  All such groups with more than two classes have abstract
algebra \(M_2(\Q)\oplus\Q\), while the multiplicity \(k(G)-2\) of the scalar
factor retains the class count at represented level.

\subsection{Cyclic groups of prime-power order}

There is a complete formula for cyclic prime powers.

\begin{lemma}
\label{lem:tree-chain-decomposition}
Let \(F\) be the linear map on \(\Q^{C_{p^a}}\) defined on coordinate
vectors by \(F\delta_x=\delta_{px}\).  There is an orthogonal rational
decomposition into reducing subspaces for \(F,F^*\) consisting of one chain
of dimension \(a+1\), together with chains of dimension \(j\) having
multiplicity
\[
 m_a=p-2,\qquad
 m_j=p^{a-j-1}(p-1)^2\quad(1\leq j<a).
 \tag{C}\label{eq:tree-chain-multiplicities}
\]
Terms of multiplicity zero are absent.  On a chain of dimension \(j\), the
algebra generated by \(F,F^*\) is \(M_j(\Q)\), and two chains of the same
length are equivalent as modules over it.
\end{lemma}

\begin{proof}
Write \(C_{p^a}=\Z/p^a\) additively and let
\(\ell(x)=a-v_p(x)\) be the number of steps from \(x\) to \(0\) under
\(x\mapsto px\), so that \(\ell(x)=j\) if and only if \(x\) has order
\(p^j\).  There are \(p^{j-1}(p-1)\) vertices at level \(j\geq1\) and one at
level \(0\).  Every vertex at level \(j\leq a-1\) has \(p\)
predecessors, all at level \(j+1\); the predecessors of \(0\) are \(0\)
itself together with the \(p-1\) elements of order \(p\).  Vertices at
level \(a\) are the leaves, and \(F^*\) annihilates them.

Note first that \(F^*\delta_x=\sum_{py=x}\delta_y\), whence
\begin{equation}
 FF^*\delta_x
 =\sum_{py=x}\delta_{py}
 =p\,\delta_x
 \qquad(\ell(x)\leq a-1),
 \label{eq:ff-star-scalar}
\end{equation}
and \(FF^*\delta_x=0\) on the leaves.  Thus \(FF^*=p\cdot I\) off the
leaves.  This is the source of the uniform chain weights below.

Nonradial chains.  Fix a vertex \(v\) with \(1\leq\ell(v)=j\leq a-1\)
and let \(w_0=\sum_yc_y\delta_y\) run over the \((p-1)\)-dimensional space
of sibling differences on the predecessor set of \(v\), so
\(\sum_yc_y=0\).  Put \(w_k=(F^*)^kw_0\) and set \(i=a-j\).  Then \(Fw_0=(\sum_yc_y)\delta_v=0\),
while \eqref{eq:ff-star-scalar} gives
\begin{equation}
 Fw_k=FF^*w_{k-1}=p\,w_{k-1}
 \qquad(1\leq k\leq i-1),
 \label{eq:chain-weights}
\end{equation}
because \(w_{k-1}\) is supported at level \(j+k\leq a-1\).  The vector
\(w_k\) is supported at level \(j+1+k\), so the chain terminates at
\(w_{i-1}\), and \(F^*w_{i-1}=0\) since that vector is
supported on leaves.  Vectors supported at distinct levels are orthogonal
for the uniform form, so the \(w_k\) are an orthogonal rational basis of a
reducing subspace of dimension \(i\).  Counting: level \(j=a-i\) carries
\(p^{a-i-1}(p-1)\) vertices for \(j\geq1\), each contributing \(p-1\)
independent differences, which gives the multiplicity
\(p^{a-i-1}(p-1)^2\) of \eqref{eq:tree-chain-multiplicities}.

Distinct chains are orthogonal.  This does not follow from level
bookkeeping alone, because two chains born at different levels do have
vectors at a common level, and their supports can even be nested.  We verify
it directly.  Since \((F^*)^k\delta_y=\sum_{z\in F^{-k}(y)}\delta_z\), the
chain vector born at \(v\) is
\begin{equation}
 w_k^{(v)}=\sum_{y\in F^{-1}(v)}c_y\,\mathbf1_{F^{-k}(y)},
 \qquad \sum_{y}c_y=0 .
 \label{eq:chain-explicit}
\end{equation}
Let \(v,v'\) be birth vertices at levels \(j\geq j'\) and let
\(w_k^{(v)},w_{k'}^{(v')}\) both be supported at level \(L\), so
\(k=L-j-1\) and \(k'=L-j'-1\), whence \(k'\geq k\).  If \(j=j'\) and
\(v\neq v'\) the fibres \(F^{-1}(v)\) and \(F^{-1}(v')\) are disjoint, hence
so are the sets \(F^{-k}(y)\) occurring in \eqref{eq:chain-explicit}, and the
inner product vanishes.  If \(j>j'\), then \(k'>k\), and for
\(y\in F^{-1}(v)\), \(y'\in F^{-1}(v')\) the intersection
\(F^{-k}(y)\cap F^{-k'}(y')\) is nonempty only if
\(y'=F^{k'-k}(y)\), in which case \(F^{-k}(y)\subseteq F^{-k'}(y')\) and the
intersection has \(p^{k}\) elements.  But every \(y\in F^{-1}(v)\) satisfies
\(F(y)=v\), so for \(k'-k\geq1\) the value
\(F^{k'-k}(y)=F^{k'-k-1}(v)=:y_0'\) is the same vertex for all such
\(y\).  Therefore
\[
 \langle w_k^{(v)},w_{k'}^{(v')}\rangle
 =p^{k}\,c'_{y_0'}\sum_{y\in F^{-1}(v)}c_y=0 .
\]
The same computation against the radial vectors \(r_L=\mathbf1_{\ell^{-1}(L)}\)
gives \(\langle w_k^{(v)},r_L\rangle=p^{k}\sum_yc_y=0\).  Hence all the chains
listed here and below are mutually orthogonal.

The radial chain.  Put \(r_k=\sum_{\ell(x)=k}\delta_x\) for
\(0\leq k\leq a\).  Then
\begin{equation}
 \begin{aligned}
 Fr_0&=r_0, & Fr_1&=(p-1)r_0, & Fr_k&=p\,r_{k-1}\ (k\geq2),\\
 F^*r_0&=r_0+r_1, & F^*r_k&=r_{k+1}\ (1\leq k\leq a-1), & F^*r_a&=0.
 \end{aligned}
 \label{eq:radial-relations}
\end{equation}
These \(a+1\) orthogonal vectors span a reducing subspace.  The predecessor
set of \(0\) contains \(0\), so \(F\) has the eigenvector \(r_0\) and is
not nilpotent here; this is the one chain on which the endpoint
argument for a nilpotent shift is unavailable.

Chains born at the origin.  The vertex \(0\) is the one birth vertex
whose predecessor set is not contained in a single level: \(F^{-1}(0)\)
consists of \(0\) itself together with the \(p-1\) elements of order \(p\).
Among the \((p-1)\)-dimensional space of sibling differences on \(F^{-1}(0)\),
those not involving \(\delta_0\), that is, supported on the level-one
vertices alone, form a subspace of dimension \((p-1)-1=p-2\).  For such a
\(w_0\) we again have \(Fw_0=(\sum_yc_y)\delta_0=0\), and \(w_0\) is supported
at level one, so the chain \(w_k=(F^*)^kw_0\) terminates at \(w_{a-1}\) and has
dimension \(a\).  This gives the multiplicity \(m_a=p-2\), absent when
\(p=2\).  The one remaining direction on \(F^{-1}(0)\), the one involving
\(\delta_0\), is already accounted for by the radial chain, as the dimension
count in \cref{thm:cyclic-prime-powers} confirms.  Note that \(i=a\) does not
arise from the nonradial analysis above, which required \(1\leq j\leq a-1\), so
the length-\(a\) chains are these \(p-2\) chains.

The algebra restricted to each chain is a full matrix algebra.  Let \(W\) be any of the
above chains, with ordered spanning vectors \(w_0,\ldots,w_{i-1}\) in the
nonradial case or \(r_0,\ldots,r_a\) in the radial case, and suppose
\(S\in\End_\Q(W)\) commutes with \(F|_W\) and \(F^*|_W\).  By
\eqref{eq:chain-weights} and \eqref{eq:radial-relations}, \(F^*|_W\) has
one-dimensional kernel, spanned by the top vector \(w_{i-1}\), respectively
\(r_a\).  Hence \(S\) preserves that line and acts on it by some
\(\lambda\in\Q\).  Applying \(F\) repeatedly and using that every weight in
\eqref{eq:chain-weights} and \eqref{eq:radial-relations} is nonzero, the
weights being \(p\), and \(p-1\geq1\) at the last radial step, gives
\(S=\lambda\) on each successive vector in turn.  Therefore
\((\PE|_W)'=\Q\), and
\cref{thm:rational-bicommutant-structural} gives
\[
 \PE|_W=(\PE|_W)''=\End_\Q(W)\cong M_{\dim W}(\Q).
\]
This argument is uniform in the two cases and does not use nilpotency.

Finally, two chains of equal length are isomorphic as modules.  For the
nonradial chains, including those born at the origin,
\eqref{eq:chain-weights} shows that \(F\) acts as \(p\) times the downward
shift and \(F^*\) as the upward shift, with the same weight \(p\) at every
step, independently of length or birth vertex; so any two such chains of the
same length are carried to one another by the linear map matching their
spanning vectors in order.  The radial chain is not of this form, since its
last step has weight \(p-1\) and \(F\) is not nilpotent on it, but it has
dimension \(a+1\), whereas every nonradial chain has dimension at most \(a\),
so no comparison between the two types is required.
\end{proof}

\begin{theorem}
\label{thm:cyclic-prime-powers}
Let \(a\geq1\).  Then
\begin{align}
 \Pram(C_{p^a})
 &\cong
 M_{a+1}(\Q)\oplus\bigoplus_{j=1}^{a}M_j(\Q),
 &&p\text{ odd},                                      \label{eq:cyclic-odd}\\
 \Pram(C_{2^a})
 &\cong
 M_{a+1}(\Q)\oplus\bigoplus_{j=1}^{a-1}M_j(\Q).
 &&                                                   \label{eq:cyclic-two}
\end{align}
Empty sums are omitted.  The \(M_{a+1}\)-module occurs once.  For odd \(p\),
the multiplicity of the \(M_a\)-module is \(p-2\), and for \(1\leq j<a\)
the multiplicity of the \(M_j\)-module is
\[
 m_j=p^{a-j-1}(p-1)^2.
\]
For \(p=2\), the \(M_a\)-block is absent and the same formula gives
\(m_j=2^{a-j-1}\) for \(j<a\).
\end{theorem}

\begin{proof}
Write \(C_{p^a}\) additively.  Its class measure is uniform, so the sole
canonical generator is pullback \(P\) along \(x\mapsto px\), together with
the ordinary transpose \(P^{\mathsf T}\).  The transpose is the map \(F\)
of \cref{lem:tree-chain-decomposition}.  That lemma gives the displayed
simple factors and their multiplicities.  The dimension identity
\[
 p^a=(a+1)+a(p-2)+
 \sum_{j=1}^{a-1}j\,p^{a-j-1}(p-1)^2
\]
shows that the listed chains exhaust the coordinate space.
\end{proof}

\begin{corollary}
\label{cor:cyclic-prime-power-dimensions}
For odd \(p\),
\[
 \dim_\Q\Pram(C_{p^a})
 =(a+1)^2+\sum_{j=1}^{a}j^2,
\]
whereas for \(p=2\) the final summand \(a^2\) is omitted.  These dimensions
depend only on \(a\) and on whether \(p=2\), while the represented
multiplicities retain \(p\).
\end{corollary}

\subsection{Dihedral and quaternion groups}
\label{sec:family-ladder}

Use the presentations
\[
 D_{2n}=\langle a,b:a^{n}=b^2=1,\ bab^{-1}=a^{-1}\rangle,
 \qquad
 Q_{4m}=\langle a,b:a^{2m}=1,\ b^2=a^m,\ bab^{-1}=a^{-1}\rangle,
\]
with all prime-power maps dividing the exponent.  Modular row reduction
followed by rational reconstruction gives the following table.

\begin{proposition}
\label{prop:family-ladder-table}
For \(3\leq n\leq16\) and \(2\leq m\leq8\), the represented algebras
\(A=\Pram(G)\) have the signatures
\[
\begin{array}{lrr|rrr@{\qquad\qquad}lrr|rrr}
G&|G|&k&\dim A&\dim Z&\dim A'&G&|G|&k&\dim A&\dim Z&\dim A'\\
\cline{1-6}\cline{7-12}
D_6   & 6&3& 9&1&1  & D_{22}&22& 7&13&5&5\\
D_8   & 8&5&17&2&2  & D_{24}&24& 9&51&3&3\\
D_{10}&10&4&10&2&2  & D_{26}&26& 8&14&6&6\\
D_{12}&12&6&26&2&2  & D_{28}&28&10&34&4&4\\
D_{14}&14&5&11&3&3  & D_{30}&30& 9&31&4&4\\
D_{16}&16&7&26&2&5  & D_{32}&32&11&41&3&11\\
D_{18}&18&6&18&3&3  &&&&&&\\
\cline{1-6}\cline{7-12}
Q_8   & 8&5&10&2&5  & Q_{24}&24& 9&51&3&3\\
Q_{12}&12&6&26&2&2  & Q_{28}&28&10&34&4&4\\
Q_{16}&16&7&26&2&5  & Q_{32}&32&11&41&3&11\\
Q_{20}&20&8&30&3&3  &&&&&&\\
\end{array}
\]
Multiplicity freeness holds at every rung except
\(D_{16}\), \(D_{32}\), \(Q_8\), \(Q_{16}\), and \(Q_{32}\), where
\(\dim A'>\dim Z\).
\end{proposition}

\begin{proof}
Rational computation as described above; the multiplicity
statement is \cref{cor:multiplicity-free} applied to the commutant and
center columns.
\end{proof}

\begin{theorem}
\label{thm:dihedral-prime-law}
Let \(p\) be an odd prime, put \(m=(p-1)/2\), and let
\[
 d=\operatorname{ord}_{U_p/\{\pm1\}}(2),\qquad q=m/d.
\]
Then
\[
 \Pram(D_{2p})\otimes_\Q\C
 \cong M_3(\C)\oplus\C^{\,d-1+\epsilon_q},
 \qquad
 \epsilon_q=
 \begin{cases}
 0,&q=1,\\
 1,&q>1.
 \end{cases}
\]
Hence
\[
 \dim_\Q\Pram(D_{2p})=8+d+\epsilon_q.
\]
The class-coordinate representation is multiplicity free if and only if
\(q=1\).
\end{theorem}

\begin{proof}
Let \(R\) be the span of the \(m\) rotation-class coordinates and let
\(\mathbf1_R\) be their sum.  The space
\[
 W=\Span_\C\{\delta_1,\delta_s,\mathbf1_R\}
\]
reduces \(\Psi_2,\Psi_p\) and their class-form adjoints.  Their restrictions
to \(W\), computed with class weights \(1,p,p-1\), generate \(M_3(\C)\).

On \(W^\perp\), the operators \(\Psi_p\) and \(\Psi_p^*\) vanish, while
\(\Psi_2\) is the permutation induced by multiplication by \(2\) on
\((\Z/p\Z)^\times/\{\pm1\}\).  This permutation has \(q\) cycles of length
\(d\).  Its eigenvalues are the \(d\)-th roots of unity.  Removing
\(\mathbf1_R\) removes the eigenvalue \(1\) when \(q=1\); for \(q>1\) its
multiplicity becomes \(q-1\).  Thus the algebra on \(W^\perp\) is
\(\C^{d-1+\epsilon_q}\).  The commutator ideal supplies the \(M_3(\C)\)
summand, so the two restrictions form the displayed direct sum.

For \(q>1\), each nontrivial \(d\)-th root occurs with multiplicity \(q\);
for \(q=1\), all eigenspaces on \(W^\perp\) are one-dimensional.  This gives
the last assertion.
\end{proof}

For \(p=17\), \(d=4\) and \(q=2\), so
\[
 \Pram(D_{34})\otimes_\Q\C\cong M_3(\C)\oplus\C^4,
 \qquad \dim_\Q\Pram(D_{34})=13.
\]

\subsection{Rational classes in \texorpdfstring{\(\mathrm{PSL}(2,q)\)}{PSL(2,q)}}

\begin{theorem}
\label{thm:psl2-rational-classes}
Let \(q=p^f\), let \(G=\mathrm{PSL}(2,q)\), and put
\[
 m_\pm=\frac{q\pm1}{\gcd(2,q-1)} .
\]
\begin{enumerate}
\item Let \(x\in G\) have order \(d\) with \(p\nmid d\).  Then \(d\) divides
\(m_+\) or \(m_-\), and the rational class of \([x]\) has \(\varphi(d)/2\)
elements if \(d>2\) and one element if \(d\leq2\).
\item The rational class of a class of \(p\)-elements has at most two
elements.
\item Let \(r\) be an odd prime.  Then \(G\) has a rational class of size
\(r\) if and only if \(\varphi(d)=2r\) for some \(d>2\) dividing \(m_+\) or
\(m_-\).  Such a \(q\) exists if and only if \(2r+1\) is prime.
\end{enumerate}
\end{theorem}

\begin{proof}
(1)  A nontrivial semisimple element of \(G\) lies in a maximal torus, and the
maximal tori of \(G\) are cyclic of order \(m_+\) or \(m_-\); this gives the
divisibility.  Let \(d>2\), let \(T\) be a maximal torus containing \(x\), and
let \(u\) be a good label, so that \(x^u\in T\) also has order \(d\).  For
\(d>2\) we have \(C_G(x)=C_G(x^u)=T\), so if \(gx^ug^{-1}=x\) then
\(gTg^{-1}=T\), that is, \(g\in N_G(T)\).  Since
\(N_G(T)/T\) has order two and acts on \(T\) by inversion, \(x^u\) is
conjugate to \(x\) if and only if \(u\equiv\pm1\pmod d\).  Thus the stabilizer
of \([x]\) in \(U_N\) is the preimage of \(\{\pm1\}\subseteq(\Z/d)^\times\),
of index \(\varphi(d)/2\).  For \(d\leq2\) every good label acts trivially on
\(\langle x\rangle\).

(2)  If \(q\) is even the \(p\)-elements are the involutions and (1) applies.
If \(q\) is odd, every nontrivial unipotent element of \(\mathrm{SL}(2,q)\)
has order \(p\) and there are two unipotent classes, so \(G\) has two
classes of \(p\)-elements in all.

(3)  By (1) and (2) the sizes of the rational classes are the numbers
\(\varphi(d)/2\) for \(d>2\) dividing \(m_+\) or \(m_-\), together with
values at most two, and \(r\geq3\); this is the stated criterion.  If
\(\varphi(d)=2r\) then \(d\) is \(\ell\) or \(2\ell\) with \(\ell=2r+1\)
prime, except for \(r=3\), where also \(d\in\{9,18\}\); in every case \(2r+1\)
is prime.  Conversely if \(\ell=2r+1\) is prime, choose by Dirichlet a prime
\(q\equiv1\pmod{2\ell}\); then \(\ell\) divides \(m_-=(q-1)/2\) and
\(\varphi(\ell)=2r\).
\end{proof}

\begin{corollary}
\label{cor:psl2-no-seven}
No \(\mathrm{PSL}(2,q)\) has a rational class of size \(7\).
\end{corollary}

\begin{proof}
By \cref{thm:psl2-rational-classes}(3) it suffices that \(14\) is not a value
of \(\varphi\).  If \(\varphi(d)=14\) then \(\varphi(p^a)\mid14\) for every
prime power \(p^a\|d\), so \(\varphi(p^a)\in\{1,2,14\}\).  For odd \(p\),
\(\varphi(p^a)=2\) gives \(p^a=3\), while \(\varphi(p^a)=14\) forces
\(p\mid14\), and \(\varphi(7)=6\), \(\varphi(49)=42\).  Hence every factor
contributes a power of two and \(\varphi(d)\) is a power of two.
\end{proof}

For \(3\leq q\leq32\) the sizes \(\varphi(d)/2\) give odd prime values only
for
\[
 \begin{array}{lll}
 \mathrm{PSL}(2,8):\ d=7,9,\ r=3, &
 \mathrm{PSL}(2,13):\ d=7,\ r=3, &
 \mathrm{PSL}(2,17):\ d=9,\ r=3,\\
 \mathrm{PSL}(2,19):\ d=9,\ r=3, &
 \mathrm{PSL}(2,27):\ d=7,14,\ r=3, &
 \mathrm{PSL}(2,23):\ d=11,\ r=5,\\
 \mathrm{PSL}(2,29):\ d=7,14,\ r=3, &
 \mathrm{PSL}(2,32):\ d=11,\ r=5, &
 \end{array}
\]
and for no other \(q\) in that range; in particular
\(\mathrm{PSL}(2,7),\mathrm{PSL}(2,9),\mathrm{PSL}(2,11),\mathrm{PSL}(2,16),
\mathrm{PSL}(2,25)\) have none.

\subsection{A comparison panel}
\label{sec:census}

\begin{theorem}
\label{thm:comparison-panel}
For each of the following 31 groups the rational algebra
\(A=\Pram(G)\) has the stated dimension, center dimension \(z\), and
represented commutant dimension \(c\); the pairs \((d,m)\) give the complex
blocks \(M_d(\C)\) and their multiplicities in \(V_G\otimes_\Q\C\).

\begingroup
\footnotesize
\normalfont
\setlength{\tabcolsep}{5pt}
\begin{longtable}{@{}lrrrrrrl@{}}
\toprule
$G$ & $|G|$ & $k(G)$ & $\exp G$ & $\dim A$ & $z$ & $c$ & $(d,m)$\\
\midrule
\endfirsthead
\toprule
$G$ & $|G|$ & $k(G)$ & $\exp G$ & $\dim A$ & $z$ & $c$ & $(d,m)$\\
\midrule
\endhead
\multicolumn{8}{@{}l}{\itshape Cyclic}\\
$C_2$ & 2 & 2 & 2 & 4 & 1 & 1 & $(2,1)$\\
$C_3$ & 3 & 3 & 3 & 5 & 2 & 2 & $(2,1),(1,1)$\\
$C_4$ & 4 & 4 & 4 & 10 & 2 & 2 & $(3,1),(1,1)$\\
$C_5$ & 5 & 5 & 5 & 5 & 2 & 10 & $(2,1),(1,3)$\\
$C_6$ & 6 & 6 & 6 & 20 & 2 & 2 & $(4,1),(2,1)$\\
$C_8$ & 8 & 8 & 8 & 21 & 3 & 6 & $(4,1),(2,1),(1,2)$\\
$C_{12}$ & 12 & 12 & 12 & 50 & 4 & 4 & $(6,1),(3,1),(2,1),(1,1)$\\
\addlinespace
\multicolumn{8}{@{}l}{\itshape Elementary abelian}\\
$C_2\times C_2$ & 4 & 4 & 2 & 5 & 2 & 5 & $(2,1),(1,2)$\\
$C_2^3$ & 8 & 8 & 2 & 5 & 2 & 37 & $(2,1),(1,6)$\\
$C_3\times C_3$ & 9 & 9 & 3 & 5 & 2 & 50 & $(2,1),(1,7)$\\
\addlinespace
\multicolumn{8}{@{}l}{\itshape Dihedral}\\
$D_8$ & 8 & 5 & 4 & 17 & 2 & 2 & $(4,1),(1,1)$\\
$D_{10}$ & 10 & 4 & 10 & 10 & 2 & 2 & $(3,1),(1,1)$\\
$D_{12}$ & 12 & 6 & 6 & 26 & 2 & 2 & $(5,1),(1,1)$\\
$D_{16}$ & 16 & 7 & 8 & 26 & 2 & 5 & $(5,1),(1,2)$\\
\addlinespace
\multicolumn{8}{@{}l}{\itshape Quaternion}\\
$Q_8$ & 8 & 5 & 4 & 10 & 2 & 5 & $(3,1),(1,2)$\\
\addlinespace
\multicolumn{8}{@{}l}{\itshape Symmetric}\\
$S_3$ & 6 & 3 & 6 & 9 & 1 & 1 & $(3,1)$\\
$S_4$ & 24 & 5 & 12 & 25 & 1 & 1 & $(5,1)$\\
$S_5$ & 120 & 7 & 60 & 49 & 1 & 1 & $(7,1)$\\
$S_6$ & 720 & 11 & 60 & 59 & 3 & 3 & $(7,1),(3,1),(1,1)$\\
$S_7$ & 5040 & 15 & 420 & 225 & 1 & 1 & $(15,1)$\\
$S_8$ & 40320 & 22 & 840 & 442 & 2 & 2 & $(21,1),(1,1)$\\
\addlinespace
\multicolumn{8}{@{}l}{\itshape Alternating}\\
$A_4$ & 12 & 4 & 6 & 10 & 2 & 2 & $(3,1),(1,1)$\\
$A_5$ & 60 & 5 & 30 & 17 & 2 & 2 & $(4,1),(1,1)$\\
$A_6$ & 360 & 7 & 60 & 27 & 3 & 3 & $(5,1),(1,1),(1,1)$\\
$A_7$ & 2520 & 9 & 420 & 65 & 2 & 2 & $(8,1),(1,1)$\\
$A_8$ & 20160 & 14 & 420 & 146 & 3 & 3 & $(12,1),(1,1),(1,1)$\\
\addlinespace
\multicolumn{8}{@{}l}{\itshape Simple of Lie type}\\
$\mathrm{PSL}(2,7)$ & 168 & 6 & 84 & 26 & 2 & 2 & $(5,1),(1,1)$\\
$\mathrm{PSL}(2,8)$ & 504 & 9 & 126 & 29 & 5 & 5 & $(5,1),(1,1),(1,1),(1,1),(1,1)$\\
$\mathrm{PSL}(2,11)$ & 660 & 8 & 330 & 38 & 3 & 3 & $(6,1),(1,1),(1,1)$\\
\addlinespace
\multicolumn{8}{@{}l}{\itshape Nonabelian of order 27}\\
$H_{27}$ & 27 & 11 & 3 & 5 & 2 & 82 & $(2,1),(1,9)$\\
\addlinespace
\multicolumn{8}{@{}l}{\itshape Frobenius of order 20}\\
$F_{20}$ & 20 & 5 & 20 & 17 & 2 & 2 & $(4,1),(1,1)$\\
\bottomrule
\end{longtable}
\endgroup

\end{theorem}

\begin{proof}
Rational row reduction gives the dimensions of the algebra, center, and
commutant, as in \cref{sec:certification}.  It also enumerates the positive
integral solutions of
\begin{equation}
 \sum_j d_j^2=\dim A,\qquad
 \sum_j m_j^2=\dim A',\qquad
 \sum_jd_jm_j=k(G),\qquad
 \#\{j\}=\dim Z(A),
 \label{eq:dimension-system}
\end{equation}
and for every row the listed solution is unique.  Modular receipts at
\(1{,}000{,}003\) and \(1{,}000{,}033\) reproduce every dimension.
\end{proof}

The groups \(D_8\) and \(Q_8\) both have five classes and identical character
tables, with \(\dim A=17\) and \(\dim A=10\); the invariant is therefore
strictly finer than the character table.  Among the symmetric groups
\(S_3,S_4,S_5\), and \(S_7\) give full matrix algebras, whereas \(S_6\) and
\(S_8\) do not.  The scalar block of \(Q_8\) occurs twice in
class-coordinate space, which is where its commutant differs from that of
\(D_8\).  Both \(H_{27}\) and \(C_3\times C_3\) have \(\dim A=5\), with
\(c=82\) and \(c=50\).

\section{Computational verification}
\label{sec:certification}

The computations use character-table data rather than matrix or permutation
realizations.  GAP~\cite{GAP} and CTblLib~\cite{CTblLib} supply the ordered
conjugacy classes, class sizes, exponents, and prime-power maps; GAP stores an
\(n\)th power map as a list whose \(i\)th entry gives the class position of
the \(n\)th powers of the elements of class \(i\)~\cite{GAPPowerMaps}.  For a
table with \(k\) classes the runner builds \(P_n\) and
\(D^{-1}P_n^{\mathsf T}D\) over \(\Q\), adjoins products until multiplication
by every generator adds no new rational row, and obtains the center and
represented commutant from rational commutator nullspaces; a separating
central element and its factorization give the central idempotents, whose
ideal dimensions give the split degrees over the center fields.  FLINT
performs the rational rank calculations~\cite{FLINTRationalMatrices}.

Reductions modulo \(1{,}000{,}003\) and \(1{,}000{,}033\) repeat every
dimension calculation.  The proof of each stated decomposition is the rational
word basis together with the central factorization; the modular runs are
consistency checks.  The GAP exports, source, tests, and JSON outputs accompany this paper as
ancillary files, together with a supplement recording the stopping criteria,
schema, file hashes, and reproduction commands.

The principal block of each sporadic algebra is certified as follows.  Since
every element of \(G\) has order dividing \(N=\exp(G)\), the operator
\(|G|^{-1}\Psi_N^*\Psi_N\) is the coordinate projection onto the identity
class, so \(E_{1A}\in A\).  The cyclic module \(W=A\delta_{1A}\) is computed
by exact rational orbit closure, and \cref{lem:cyclic-rank-one} gives
\(A|_W=\End_\Q(W)\); the degree of the principal block is therefore
\(\dim_\Q W\), replayed at two primes by nonzero orbit minors.  The residual
algebra on the orthogonal complement of \(W\) is then closed and decomposed
separately.  This replaces an earlier procedure that took the number of
cells of the stable weighted refinement as the principal degree.  That
partition is equitable, but equitability does not imply fullness, and for
\(Fi_{23}\) it overcounts: the refinement has \(92\) cells while
\(\dim_\Q A\delta_{1A}=91\), the missing dimension being the invariant
line of \cref{sec:invisible-symmetries}.  The other twenty-five sporadic
rows are unchanged by the corrected procedure.

Independently of the word-basis ranks, every reported row is checked against
the represented Wedderburn identities
\[
 \sum_jd_j^2=\dim A,\qquad
 \sum_jm_j^2=\dim A',\qquad
 \sum_jd_jm_j=k(G),\qquad
 \sum_j[K_j:\Q]=\dim_\Q Z(A),
\]
where conjugate complex blocks are grouped over their rational center.
These checks are especially restrictive in the largest rows.  For the
Monster they give
\[
 170+5+3+4\cdot2+8=194,\qquad
 170^2+5^2+3^2+4\cdot2^2+8=28{,}958,
\]
with fifteen multiplicity-one rational blocks.  For \(O'N\), the quadratic
center contributes two dimensions to the center and the class count forces
one scalar block of multiplicity two.  For \(\B\), the center and commutant
dimensions force two multiplicity-two blocks.  Thus the tables are
subject both to exact algebra closure and to the full set of representation
identities, not only to modular rank agreement.

The ramified local factors are computed by two routes.  The first forms the
\(I_{p,m}\)-orbits and the induced Frobenius permutation; the second uses,
for each rational-class stabilizer \(S\leq U_m\),
\[
 f=\operatorname{ord}_{U_m/SI_{p,m}}(\kappa_{p,m}),\qquad
 g=[U_m:SI_{p,m}]/f,
\]
and the factor \((1-T^f)^g\).  The two routes agree in all \(173\)
sporadic bad-prime slots.  The exported ramified maps generate the required
\(U_m\)-action in \(168\) slots; the other five are omitted from the
computed table.  Cyclic groups through order \(72\) and eight abelian
groups give the cyclotomic local factors independently.


\section{Further problems}
\label{sec:open-problems}

The inverse problem asks which semisimple \(\Q\)-algebras occur as
\(\Pram(G)\), or as \(\PE(\mathscr X)\) with the class-coordinate module
attached.  The identities \eqref{eq:dimension-system} are necessary but far
from sufficient: by \cref{thm:prime-exponent-structural} every group of prime
exponent with \(k>2\) gives \(M_2(\Q)\oplus\Q\) whatever its order.  For the
extreme case \cref{cor:full-matrix} converts the question into rigidity:
\(\dim A=k^2\) if and only if the weighted colored power graph has no
balanced kernel beyond the scalars.

By \cref{thm:psl2-rational-classes} the odd primes \(r\) for which \(2r+1\) is
composite are unrealized in \(\mathrm{PSL}(2,q)\), the first being \(r=7\).
They are realized by abelian groups: \(C_{29}\) has a single nonidentity
rational class, of size \(28\), and hence a summand with endomorphism field
\(\Q(\zeta_7)\).  Which \(\Q(\zeta_r)\) occur as centers of blocks of
\(\Pram(G)\) for \(G\) simple is open; the value \(r=7\) occurs for no
sporadic group.

The Monster, \(Co_2\), \(Co_3\), and \(Fi_{23}\) each have a nonprincipal
simple factor inside the Galois-fixed part of \(V_G\)
(\cref{sec:invisible-symmetries}).  For the Monster it comes from a
permutation of classes; for \(Fi_{23}\) it does not.  Which finite groups
admit symmetries of the weighted power system that are not automorphisms of
the character table, and which admit balanced kernels beyond the span of
all such symmetries, are open.

Equality in the synchronized tensor containment
\eqref{eq:direct-product-containment} is open.

\section*{Code and data availability}

The GAP and CTblLib exports, the rational word-closure runner, the
Monster certificate, and the JSON outputs underlying
\cref{thm:comparison-panel}, \cref{thm:sporadic-census},
\cref{thm:good-label-census}, and \cref{thm:monster-symmetry-commutant}
are included as ancillary files, with a single reproduction script and a
test suite.  \Cref{sec:certification} describes the checks they perform.

\appendix
\section{Power dynamics and the sporadic separation theorem}
\label{sec:power-dynamics}

The forward operators retain a finite dynamical invariant that is different
from the adjoint-closed algebra.  Let \(\tau:X\to X\) be a map of a finite
set and let \(P_\tau\) act on \(\Q^X\) by pullback,
\((P_\tau f)(x)=f(\tau x)\).  Write \(c_\ell(\tau)\) for the number of
directed cycles of length \(\ell\) in the functional graph of \(\tau\).

\begin{theorem}[Finite-map zeta formula]
\label{thm:finite-map-zeta}
For every finite map \(\tau:X\to X\),
\begin{equation}
 Z_\tau(z):=\det(I-zP_\tau)^{-1}
 =\prod_{\ell\geq1}(1-z^\ell)^{-c_\ell(\tau)}.
 \label{eq:finite-map-zeta}
\end{equation}
Thus \(Z_\tau\) depends only on the periodic cycles.  The rooted trees
feeding those cycles contribute no factor.
\end{theorem}

\begin{proof}
Order the vertices of each connected functional-graph component by
decreasing distance from its unique directed cycle, followed by the cycle
vertices.  In this basis \(P_\tau\) is block triangular.  Its off-cycle
diagonal block is nilpotent, while a cycle of length \(\ell\) contributes
the permutation matrix of an \(\ell\)-cycle.  Hence
\[
 \det(I-zP_\tau)
 =\prod_{\text{cycles }C}\det(I-zP_C)
 =\prod_{\text{cycles }C}(1-z^{|C|}).
\]
\end{proof}

For a finite group \(G\), put \(N_G=\exp(G)\) and
\[
 \mathcal N_G=\{p^a:p\text{ prime},\ a\geq1,\ p^a\mid N_G\}.
\]
Let \(\sigma_n:\Conj(G)\to\Conj(G)\) be \(C\mapsto C^n\).  The
\emph{ramified dynamical packet} is the unlabelled multiset
\begin{equation}
 \mathfrak Z_{\rm ram}(G)
 =\bigl\{\!\bigl\{Z_{\sigma_n}(z):n\in\mathcal N_G\bigr\}\!\bigr\}.
 \label{eq:ramified-dynamical-packet}
\end{equation}
It forgets both the prime-power labels and the ordering of the conjugacy
classes.  It is therefore weaker than the fully marked eidetic object.

\begin{theorem}[Sporadic dynamical separation]
\label{thm:sporadic-dynamical-separation}
For the twenty-six sporadic simple groups:
\begin{enumerate}
\item the map \(G\mapsto\mathfrak Z_{\rm ram}(G)\) is injective;
\item the cycle lengths occurring among all
      \(\sigma_n\), \(n\in\mathcal N_G\), lie in
      \(\{1,2,3,5\}\);
\item for each \(G\), the set of occurring cycle lengths equals
      \(\{1\}\) together with the orders of the good-label characters in
      \cref{thm:good-label-census};
\item a cycle of length \(3\) occurs exactly for
      \[
      J_1,\ J_3,\ J_4,\ Ly,\ O'N,\ Ru,
      \]
      and a cycle of length \(5\) occurs exactly for \(Ly\).
\end{enumerate}
\end{theorem}

\begin{proof}
The class power maps in the fixed CTblLib ordering give \(258\)
prime-power directions.  For each direction, cycle enumeration followed by
\cref{thm:finite-map-zeta} gives its exponent vector
\((c_1,c_2,c_3,c_5)\).  Sorting these vectors within each group produces
twenty-six distinct multisets, proving (1).  Direct enumeration gives
(2).  Comparing the cycle-order support with the independently reconstructed
good-label characters of \cref{thm:good-label-census} proves (3).
The same exact enumeration gives the two lists in (4).
\end{proof}

\begin{corollary}[Odd cyclotomic detection]
\label{cor:dynamics-cyclotomic-detection}
Among the sporadic simple groups, the ramified dynamical packet detects
every nonrational center field in \(\Pram(G)\).  Cubic cycles occur exactly
for the six groups carrying a \(\Q(\zeta_3)\)-factor, while the unique
quintic cycle occurs for the unique group carrying a
\(\Q(\zeta_5)\)-factor.
\end{corollary}

\begin{proof}
Combine \cref{thm:sporadic-dynamical-separation} with
\cref{thm:sporadic-census,thm:good-label-census}.
\end{proof}

\begin{remark}
The equality of supports in
\cref{thm:sporadic-dynamical-separation}(3) is a census theorem, not a
general equivalence for finite groups.  The general implication supplied
by \cref{lem:orbit-exponent-bound} is one-sided: character order is bounded
by permutation order on block support.
\end{remark}


\begin{thebibliography}{99}

\bibitem{Adams}
J.~F. Adams,
\newblock Vector fields on spheres,
\newblock \emph{Annals of Mathematics} \textbf{75} (1962), 603--632.
\newblock MR0139178.
\newblock \url{https://doi.org/10.2307/1970213}.

\bibitem{AtiyahTall}
M.~F. Atiyah and D.~O. Tall,
\newblock Group representations, \(\lambda\)-rings and the \(J\)-homomorphism,
\newblock \emph{Topology} \textbf{8} (1969), 253--297.
\newblock MR0244387.
\newblock \url{https://doi.org/10.1016/0040-9383(69)90015-9}.

\bibitem{BannaiIto}
E.~Bannai and T.~Ito,
\newblock \emph{Algebraic Combinatorics I: Association Schemes},
\newblock Benjamin/Cummings, 1984.
\newblock MR0882540.

\bibitem{CTblLib}
T.~Breuer,
\newblock \emph{CTblLib: The GAP Character Table Library},
\newblock version 1.3.9, 2024,
\newblock \url{https://www.math.rwth-aachen.de/~Thomas.Breuer/ctbllib/}.


\bibitem{ChenPonomarenko}
G.~Chen and I.~Ponomarenko,
\newblock \emph{Lectures on Coherent Configurations},
\newblock Central China Normal University Press, 2019.
\newblock \url{https://www.pdmi.ras.ru/~inp/ccNOTES.pdf}.

\bibitem{Atlas}
J.~H. Conway, R.~T. Curtis, S.~P. Norton, R.~A. Parker, and R.~A. Wilson,
\newblock \emph{Atlas of Finite Groups},
\newblock Oxford University Press, 1985.
\newblock MR0827219.

\bibitem{Wedderburn}
C.~W. Curtis and I.~Reiner,
\newblock \emph{Methods of Representation Theory}, Vol.~I,
\newblock Wiley, 1981.
\newblock MR0632548.

\bibitem{Davenport}
H.~Davenport,
\newblock \emph{Multiplicative Number Theory},
\newblock third edition, Springer, 2000.
\newblock MR1790423.

\bibitem{FLINTRationalMatrices}
The FLINT Development Team,
\newblock \emph{FLINT: Fast Library for Number Theory}, version 3.6.0, 2025.
\newblock Documentation: matrices over the rational numbers.
\newblock \url{https://flintlib.org/doc/fmpq_mat.html}.

\bibitem{GAP}
The GAP Group,
\newblock \emph{GAP -- Groups, Algorithms, and Programming},
\newblock version 4.14.0, 2024.
\newblock \url{https://www.gap-system.org}.

\bibitem{GAPClassFunctions}
The GAP Group,
\newblock \emph{GAP Reference Manual}, Chapter 72: Class Functions,
\newblock \url{https://gap-system.github.io/gap/doc/ref/chap72_mj.html}.

\bibitem{GAPPowerMaps}
The GAP Group,
\newblock \emph{GAP Reference Manual}, Chapter 73: Maps Concerning Character
Tables,
\newblock \url{https://gap-system.github.io/gap/doc/ref/chap73_mj.html}.

\bibitem{HigmanCoherent}
D.~G. Higman,
\newblock Coherent configurations I,
\newblock \emph{Rendiconti del Seminario Matematico della Universit\`a
di Padova} \textbf{44} (1970), 1--25.
\newblock MR0325420.
\newblock \url{https://www.numdam.org/item/RSMUP_1970__44__1_0/}.

\bibitem{HigmanCoherentAlgebras}
D.~G. Higman,
\newblock Coherent algebras,
\newblock \emph{Linear Algebra and its Applications} \textbf{93} (1987),
209--239.
\newblock MR0898557.
\newblock \url{https://doi.org/10.1016/S0024-3795(87)90326-0}.

\bibitem{Isaacs}
I.~M. Isaacs,
\newblock \emph{Character Theory of Finite Groups},
\newblock Dover, 1994.
\newblock MR1280461.

\bibitem{Knutson}
D.~Knutson,
\newblock \emph{\(\lambda\)-Rings and the Representation Theory of the
Symmetric Group},
\newblock Lecture Notes in Mathematics 308, Springer, 1973.
\newblock MR0364425.

\bibitem{QureshiReis}
C.~Qureshi and L.~Reis,
\newblock On the functional graph of the power map over finite groups,
\newblock \emph{Discrete Mathematics} \textbf{346} (2023), no.~7,
Article 113393.
\newblock MR4562433.
\newblock \url{https://doi.org/10.1016/j.disc.2023.113393}.

\bibitem{Reiner}
I.~Reiner,
\newblock \emph{Maximal Orders},
\newblock London Mathematical Society Monographs 28, Oxford University Press,
2003.
\newblock MR1972204.

\bibitem{Serre}
J.-P. Serre,
\newblock \emph{Linear Representations of Finite Groups},
\newblock Graduate Texts in Mathematics 42, Springer, 1977.
\newblock MR0450380.

\bibitem{SteinbergTransformation}
B.~Steinberg,
\newblock A theory of transformation monoids: combinatorics and
representation theory,
\newblock \emph{Electronic Journal of Combinatorics} \textbf{17}
(2010), no.~1, Research Paper 164.
\newblock MR2745707.
\newblock \url{https://doi.org/10.37236/436}.

\bibitem{SteinbergMonoids}
B.~Steinberg,
\newblock \emph{Representation Theory of Finite Monoids},
\newblock Universitext, Springer, 2016.
\newblock MR3525092.
\newblock \url{https://doi.org/10.1007/978-3-319-43932-7}.

\bibitem{Washington}
L.~C. Washington,
\newblock \emph{Introduction to Cyclotomic Fields},
\newblock second edition, Graduate Texts in Mathematics 83, Springer, 1997.
\newblock MR1421575.


\bibitem{Weisfeiler}
B.~Weisfeiler (ed.),
\newblock \emph{On Construction and Identification of Graphs},
\newblock Lecture Notes in Mathematics 558, Springer, 1976.
\newblock MR0543783.

\bibitem{Yamada}
T.~Yamada,
\newblock \emph{The Schur Subgroup of the Brauer Group},
\newblock Lecture Notes in Mathematics 397, Springer, 1974.
\newblock MR0347957.

\end{thebibliography}

\end{document}
